\input{preamble}

\begin{document}

\title{Supplement to ``From Obstruction to Rate: The Observation Order of a Safe Compression
Decision'':\\ Full Development, Formal Declarations and Status}

\author{\IEEEauthorblockN{Vu Phan An Nguyen}
\IEEEauthorblockA{\textit{Rider University} \\
Lawrenceville, NJ, USA}}

\maketitle

\begin{abstract}
This supplement carries the full development behind the main paper: the latent family and its
safe sets, the moment law, the adaptive transcript on finite alphabets and on filtrations, the
affinity route, the envelope and the density identification, the combinatorial law and its
error floor, the $\log(1/\delta)$ side, the headline assembly, the geometry and its density
realisation, the composite test and decoder, the per-pair likelihood-ratio test, the parity
channel and its decoders, and the crossover mechanism, each with its Lean declarations. Nothing
of the pre-rewrite text is dropped. Where a later theorem of the main paper supersedes a
statement or closes an item called open here, the list below says so; the body is kept as the
record, with one correction made in place (the decoder note).
\end{abstract}

\section*{Status map (2026-09-25)}
\begin{itemize}
\item \emph{The general rate law}, called open in the combinatorial section below, is now
Theorems ``The lower law'', ``The tournament'' and ``The order law'' of the main paper
(\lean{obstruction\_lower\_star}, \lean{tournament\_law}, \lean{graded\_tStar} in
\lean{LeanPaperBObstructionLawV1}).
\item \emph{The parity rates} are also corollaries of that law
(\lean{parity\_safe\_from\_law}, \lean{parity\_id\_from\_law} in
\lean{LeanPaperBObstructionParityV1}).
\item \emph{The continuous adaptive transcript}, called the open middle link below, is carried
for every channel with strictly positive measurable real densities in protocol form by
\lean{transcript\_kl\_le} and
\lean{transcript\_kl\_le\_qsum}, which the general lower law uses.
\item \emph{The latent family as a channel} (called open in the transcript and headline sections below). Its channel is now
written in protocol form over every finite family of linear probes, and its two nuisance worlds
are identified exactly with the iterated-difference pair (\lean{famChan}, \lean{famDens\_vdiff},
\lean{famDens\_udiff}, \lean{famChi\_v\_le}, \lean{famChi\_u\_le} in
\lean{LeanPaperBFamilyChannelV1}). Both adaptive lower bounds hold on it
(\lean{family\_safe\_lower}, \lean{family\_id\_lower}).
\item \emph{Upper bounds on the Gaussian model} (the Chebyshev side of the headline section below). Explicit decoders with
$\log(1/\delta)$ budgets replace the Chebyshev $1/\alpha$ side (\lean{safe\_decoder\_error},
\lean{id\_decoder\_error}, \lean{family\_safe\_minimax}, \lean{family\_id\_minimax} in
\lean{LeanPaperBFamilyDecoderV1}). The statements of \lean{LeanPaperBHeadlineV1} remain true,
and they are superseded as the rates on this family.
\item \emph{Unbounded stopping on every channel.} The truncation argument, proved below only
on the parity channel, is now proved for every channel in protocol form
(\lean{stopped\_obstruction\_lower\_star}, \lean{trunc\_correct}, \lean{correctSt\_ofFinite}
in \lean{LeanPaperBStoppedLawV1}), and on the latent family
(\lean{family\_safe\_stopped\_rate}, \lean{family\_id\_stopped\_rate}).
\item \emph{The tournament's tests on every channel.} Block likelihood-ratio tests discharge the
test hypothesis of the upper law with the Hellinger information (\lean{block\_error},
\lean{hellinger\_tournament}, \lean{hellinger\_upper\_star}, \lean{resolved\_le\_scale},
\lean{hellinger\_tournament\_satisfiable} in \lean{LeanPaperBHellingerTestV1}).
\item \emph{The general law is tight on both channels.} A bounded feature's mean gap is at most
$\int|p-q|\le\sqrt8\,H$, so the channels' own feature gaps bound their Hellinger information
below, and the general tournament attains all four rates (\lean{tvs\_sq\_le},
\lean{gap\_sq\_le\_hel}, \lean{family\_safe\_general\_upper}, \lean{family\_id\_general\_upper} in
\lean{LeanPaperBHellingerFamilyV1}; \lean{parity\_safe\_general\_upper},
\lean{parity\_id\_general\_upper} in \lean{LeanPaperBHellingerParityV1}).
\item \emph{The latent family's safe sets are derived}, not stated (\lean{famSafe\_derived} in
\lean{LeanPaperBFamilySafeV1}), for $r\ge2$.
\item \emph{Scope, not open}: the constants depend on $r$ explicitly and are not uniform in it.
\emph{Optional}: a two-sided bound on the information over the class $Q_{\eta}$.
\item \emph{The decoder brackets} are corrected in place (note after the decoder figure).
\end{itemize}

\section{The family, and its safe sets}
\label{sec:family}

For $r\ge2$ and $u,v\in\{-1,1\}$ put $X=\{-1,1\}^{r+3}$ and
\begin{equation}
F_{u,v}(z)=3u z_{1}-2z_{2}-2z_{3}+v\prod_{\ell=4}^{r+3}z_{\ell},
\label{eq:family}
\end{equation}
with the uniform product reference and the deployment support
\begin{equation}
K=\bigl\{(t,t,t,1,\dots,1):t=\pm1\bigr\}.
\label{eq:support}
\end{equation}
The single-coordinate responses are $Q_{1}(z_{1})=3uz_{1}$, $Q_{2}=-2z_{2}$, $Q_{3}=-2z_{3}$
and $Q_{j}=0$ for $j\ge4$, and the losses on $K$, computed as half the oscillation of the
retained residual, are
\begin{equation}
\begin{gathered}
L_{K}(\{1\})=4,\qquad L_{K}(\{2\})=L_{K}(\{3\})=|3u-2|,\\
L_{K}(\{j\})=|3u-4|
\end{gathered}
\label{eq:losses}
\end{equation}
for $j\ge4$. Every entry is derived from~\eqref{eq:family} and~\eqref{eq:support}; no loss
table is declared. \CERT{}

Two consequences matter. For $\varepsilon<1$ the safe sets at $u=+1$ and $u=-1$ are
\emph{disjoint}---at $u=-1$ only $\{1\}$ is safe, at $u=+1$ every other singleton is---so a
safe compression rule also decides $u$. And $v$ does not move the safe set at all: the
interaction sign is decision-irrelevant.

\subsection{The observation model, which is where the content is}

A query picks a linear $g$ with $\|g\|_{\infty}\le1$ and returns
\begin{equation}
Y=F_{u,v}(Z)+h\,g(Z)+\xi,\qquad Z\sim\mathrm{Unif}(X),\ \xi\sim N(0,\sigma^{2}).
\label{eq:obs}
\end{equation}
Three conditions are not decoration:
\begin{itemize}
\item $Z$ is \emph{latent} and is never returned;
\item queries may sample on $X$ while the compression is judged on $K$;
\item no query reads $F$ directly, and $g$ may not be amplified by $1/h$.
\end{itemize}

\TODOPAR{If $Z$ were observed, or the class were richer than bounded linear probes, the
lower bound can collapse. If an application forces measurement on $K$ only, the construction
must change; that difference must not be hidden in notation. State this beside the model in
the final version, not in a remark.}

A fact worth isolating: $u$ is \emph{invisible} at order zero. The law of $F$ alone does not
depend on $\mathrm{sign}(u)$, because $z_{1}$ is symmetric; only the probe $h\,g$ breaks that
symmetry. This is why the decision costs $h^{-2}$ rather than $O(1)$, and it is a property of
the construction rather than an assumption. \CERT{}

\section{The order law}
\label{sec:order}

\subsection{Vanishing}

\begin{theorem}[Moment vanishing]
\label{thm:vanish}
Let $\chi$ be the top character on $\{-1,1\}^{n}$ and let $g$ be any polynomial of degree at
most $d$. Then $\mathbb{E}[\chi\,g^{j}]=0$ for every $j$ with $jd<n$, in particular for every
$j<\lceil n/d\rceil$.
\end{theorem}

\leanbox{\lean{moment\_vanishes\_below\_order}, \lean{moment\_vanishes\_below\_s}, with
\lean{chiS\_mul}, \lean{HasDeg\_mul}, \lean{top\_character\_orthogonal}, in
\lean{LeanUnifiedQueryOrderV1}}

The proof is short and is the reason the statement is about \emph{every} admissible query
rather than a chosen one: characters multiply to the character of the symmetric difference,
so degrees add under multiplication, and the only set surviving against the top character is
the full one, which needs $jd\ge n$.

\subsection{Attainment, with the exact constant}

\begin{theorem}[Attainment]
\label{thm:attain}
For the averaged tail query,
$\displaystyle\sum_{z\in\{-1,1\}^{n}}\Bigl(\prod_{i}z_{i}\Bigr)\Bigl(\sum_{i}z_{i}\Bigr)^{k}$
vanishes for $k<n$ and equals $2^{n}\,n!$ for $k=n$. Normalising,
$\mathbb{E}[\chi_{n}g^{n}]=n!/n^{n}\neq0$.
\end{theorem}

\leanbox{\lean{cubeSum\_below\_diag}, \lean{cubeSum\_diag}, \lean{tail\_query\_moment},
\lean{tail\_query\_moment\_ne\_zero} in \lean{LeanPaperBAttainmentV1}}

The proof is a recurrence, not a bijection count: peeling one coordinate and expanding the
binomial gives
\begin{equation}
\begin{split}
&\mathrm{cubeSum}(n{+}1,k)\\
&\quad=\sum_{j}\binom{k}{j}\bigl(1+(-1)^{j+1}\bigr)\,\mathrm{cubeSum}(n,k{-}j),
\end{split}
\label{eq:rec}
\end{equation}
whose parity factor kills every even $j$, so the recursive index always drops strictly and a
single induction gives both lines. Theorems~\ref{thm:vanish} and~\ref{thm:attain} together
turn an $\Omega$ into a $\Theta$.

\subsection{Both scales, on the real family}

On~\eqref{eq:family}, the decision probe $g=z_{1}$ gives
$\mathbb{E}[Y^{2}]=18+\sigma^{2}+6uh+h^{2}$ exactly, so the two values of $u$ sit $12h$
apart; and the order-$(r{+}1)$ moment contrast between $v=\pm1$ under the tail query is
exactly $2(r{+}1)\,r!/r^{r}\,h^{r}$---an identity, not a leading term. \CERT{}

Measuring the squared Hellinger distance of the \emph{actual} observation laws by
deterministic quadrature gives log--log slopes $1.983$ for the decision pair against a
predicted $2$, and $3.959$ and $5.962$ for the identification pair at $r=2,3$ against
predicted $4$ and $6$. \NUM{} Under the wrong probe the identification pair is not merely
higher-order but \emph{exactly} indistinguishable: $H^{2}<10^{-12}$ at every $h$, because
flipping $v$ is a measure-preserving relabelling of the tail coordinates and $z_{1}$ is
independent of them. The order is a property of the query class, not only of the worlds.

\section{Against every adaptive policy}
\label{sec:adaptive}

A per-query bound says nothing about a learner that chooses its next query from what it has
seen. The step that makes the lower bound legitimate is the transcript chain rule, and we
prove it for a transcript law built from the policy rather than assumed to be a product.

\begin{theorem}[Transcript chain rule and horizon bound]
\label{thm:chain}
For a policy $\mathrm{pol}$ mapping histories to queries,
\begin{equation}
\mathrm{KL}_{n+1}(h)=\mathrm{kl}\bigl(P_{\mathrm{pol}(h)},P'_{\mathrm{pol}(h)}\bigr)
+\sum_{o}P_{\mathrm{pol}(h)}(o)\,\mathrm{KL}_{n}(h\!\cdot\!o),
\label{eq:chain}
\end{equation}
and hence $\mathrm{KL}_{N}\le N\sup_{q}\mathrm{kl}(P_{q},P'_{q})$.
\end{theorem}

\leanbox{\lean{klPath\_chain}, \lean{klPath\_le\_horizon}, with \lean{sumPaths\_mass} and
\lean{klPath\_nonneg}, in \lean{LeanPaperBTranscriptKLV1}}

\paragraph{The same bound without a finite alphabet.}
The chain rule above sums over every length-$n$ list of observations, so it needs a finite
observation alphabet --- and the construction's observation is Gaussian. The object that
carries the argument without that assumption is the likelihood ratio on the filtration,
$L_{t}=\mathrm{d}P|_{\mathcal F_{t}}/\mathrm{d}Q|_{\mathcal F_{t}}$. Four facts, none of which
needs the alphabet to be finite: $\int f\,\mathrm{d}P=\int L_{t}f\,\mathrm{d}Q$ for every
$\mathcal F_{t}$-measurable $f$; $\int_{A}L_{t+1}\,\mathrm{d}Q=\int_{A}L_{t}\,\mathrm{d}Q$ for
every $A\in\mathcal F_{t}$, because both compute $P(A)$; that consistency \emph{is} the tower
property $L_{t}=\mathbb{E}_{Q}[L_{t+1}\mid\mathcal F_{t}]$; and with $L_{0}=1$ the divergence
telescopes, $\mathbb{E}_{P}\log L_{n}=\sum_{t<n}\mathbb{E}_{P}\log(L_{t+1}/L_{t})$, giving the
horizon bound $nB$ under a per-step bound $B$. The policy never appears: it is inside the
filtration, so a policy reading the entire past is already covered --- a strengthening over
the list formulation, where cancelling the policy is a step of the proof.

\NOTEPAR{What is still owed is the bridge, not the analysis: that the transcript law of a
concrete adaptive Gaussian policy admits such an $L_{t}$, built from that policy's own
kernels. Until that is written the continuous adaptive transcript is not carried, and this
paragraph must not be read as saying it is. Integrability of each $\log L_{t}$ under $P$ is a
hypothesis.}

\leanbox{\lean{change\_of\_measure}, \lean{setIntegral\_L\_succ}, \lean{condExp\_L},
\lean{kl\_eq\_sum\_increments}, \lean{kl\_le\_horizon} in \lean{LeanPaperBChangeV1}}

The policy contributes nothing because it is the same map under both worlds. With the
per-query bound of Theorem~\ref{thm:vanish} this gives $N\gtrsim h^{-2r}$ for identification
against every adaptive policy \emph{on a finite observation alphabet}. The observation model
of this paper has Gaussian noise, so Theorem~\ref{thm:chain} does not close the lower bound
on the model we actually use, and we do not claim that it does.

\subsection{What would close it, and what is already proved of that}

The route that does close it runs on the Hellinger \emph{affinity}
$\rho(P,Q)=\int\sqrt{pq}=1-H^{2}(P,Q)$, because a policy that uses the same query kernel
under both worlds contributes a factor of one at every step, so
\begin{equation}
\rho\bigl(\mathbb{P}^{\pi,N}_{+},\mathbb{P}^{\pi,N}_{-}\bigr)\ \ge\
\bigl(1-C_{r}h^{2r}\bigr)^{N},
\label{eq:affinity}
\end{equation}
\emph{provided} the per-step bound $\sup_{g}H^{2}(P^{g}_{+,h},P^{g}_{-,h})\le C_{r}h^{2r}$
holds uniformly over the query class. Two-point testing then gives
$\mathrm{err}_{N}\ge\tfrac14\rho^{2}$.

Everything after that per-step bound is arithmetic, and it is proved: the product in
\eqref{eq:affinity} from a per-step floor, Bernoulli's inequality, and the conclusion that no
budget below $(1-2\sqrt{\delta_{0}})/(C_{r}h^{2r})$ reaches error $\delta_{0}$. The same
statement with $h^{2}$ in place of $h^{2r}$ gives the decision side. So the general adaptive
lower bound is reduced to \emph{one} analytic obligation---a per-step Hellinger bound
uniform over the query class---and to nothing else.

\leanbox{\lean{affinity\_pow\_lower}, \lean{affinity\_bernoulli},
\lean{budget\_lower\_from\_affinity}, \lean{two\_exponents\_from\_affinity} in
\lean{LeanPaperBRatesV1}}

\paragraph{Where this sits, and what it is not.}
The \emph{phenomenon} behind \eqref{eq:affinity} has a clear precedent. The masking and
side-channel literature has studied exactly this mechanism for years: low moments carry no
information, the signal appears only at high order, a coupling or a nonlinearity lowers
the protective order, and changing the noise level changes which regime
dominates~\cite{duc2015}. Duc, Faust and Standaert discuss non-independent leakage
directly. Their treatment has heuristic components and does not contain the adaptive
theorem proved here --- but it does mean that
\[
H^{2}\asymp h^{2r},\qquad N\asymp h^{-2r},
\]
and even the $h^{6}+\eta^{2}h^{4}$ crossover, are \textbf{not} strong enough on their own
to carry a paper. They are recorded here as supporting results.

Likewise, keeping only decision-relevant information is itself an established
idea~\cite{poli2023}. What would have to be new is the \emph{exponent law and its
geometric characterisation}, not the slogan.

The question we therefore keep is:

\begin{quote}
\emph{the order of information needed to find a safe compression can differ from the order
needed to identify the world, already when the safe answer sets overlap and the queries are
chosen adaptively.}
\end{quote}

with the target
\begin{gather*}
N_{\rm safe}(h,\delta)\asymp h^{-2s_{\rm safe}}\log(1/\delta),\\
N_{\rm id}(h,\delta)\asymp h^{-2s_{\rm id}}\log(1/\delta),
\qquad s_{\rm safe}<s_{\rm id},
\end{gather*}
and three conditions that make it worth more than an easy-target-plus-hard-nuisance
construction: the safe sets must come from the actual \lean{selectedRadius} geometry rather
than an assigned answer table; they must genuinely \emph{overlap}, as in
$\{1,2\},\{2,3\},\{1,3\}$, so that no single nuisance bit can be discarded; and the lower
bound must cover every adaptive policy with a matching upper bound on the same query class.

\NOTEPAR{None of those three is delivered here. What is delivered is the analytic half:
\eqref{eq:affinity} on the linear class with explicit constants, and the
$\log(1/\delta)$ side of the budget below. If the geometry half does not follow, this
material belongs inside the Support paper as a supporting rate, not as a paper of its own.
That is the honest reading and we do not argue against it. \emph{Status 2026-09-26: superseded. All three are now delivered: safe sets computed from the
loss (\lean{safe\_iff\_ne}), the overlapping correspondence (\lean{coSafe\_pairwise\_meet}), and
matched adaptive rates on one query class (main text, Theorems ``The two rates'' and ``Both
rates on the latent family'').}}

\paragraph{The remaining obligation, measured.}
Since everything after \eqref{eq:affinity} is proved, the whole adaptive lower bound turns
on one number, and that number can be computed. Both densities are exact finite Gaussian
mixtures, so $\sup_{g}H^{2}$ is a quadrature over the mixtures themselves, with the
supremum attacked by many restarts of a projected ascent over the probe class, with the
canonical construction probe $(z_{4}+z_{5}+z_{6})/3$ seeded as a mandatory baseline rather
than left for the random restarts to rediscover. On the \emph{linear} class the fitted
identification orders are $3.79$ at $r=2$ and $5.80$ at $r=3$ against the claimed $2r=4$
and $6$, and $1.91$ in the decision direction against the claimed $2$. Fitting $C_{r}$ on
the coarse half and testing it on the fine half gives worst ratios $1.089$ at $r=2$ and
$0.126$ at $r=3$. That test is an \emph{upper-envelope} check and nothing more: a ratio
well below one means the measured order \emph{exceeds} $2r$, so the envelope holds with
room; it is not a two-sided approximation and is not reported as one.

\NOTEPAR{Two implementation faults produced an earlier $r=3$ slope of $8.13$, and both are
repaired. The squared Hellinger distance was formed as $1-\rho$ with $\rho$ the affinity;
at $H^{2}\sim10^{-8}$ that subtracts two numbers agreeing to eight digits and discards half
the double-precision mantissa before quadrature error is even counted. It is now computed as
$\tfrac12\int(p-q)^{2}/(\sqrt p+\sqrt q)^{2}$, which has no such cancellation. And the probe
seeded as ``canonical'' was a flat vector over every coefficient, which spreads the same mass
over the head coordinates $z_{1},z_{2},z_{3}$; the construction's attaining probe is the
tail average $(1/r)\sum_{\ell\ge4}z_{\ell}$ alone, now recovered from the basis by least
squares with the residual checked ($6\cdot10^{-16}$) and recorded. With both repaired the
slope is $5.80$. \emph{This is better evidence, not a proof:} the ascent still fails to
improve on its best seed at every $h$, so the value remains a lower bound on the supremum
and must not be read as evidence for an $h^{6}$ optimum law. The artifact records this per $h$ as
\texttt{sup\_equals\_best\_random} and \texttt{n\_search\_limited}. What
\eqref{eq:affinity} consumes is only the domination by $C_{r}h^{2r}$ over the measured
range, and a value attained by a legal probe is a \emph{lower} bound on the supremum ---
the direction that can refute an upper-bound law and can never certify one.}

\paragraph{The query class must be the linear one, and that is a theorem-level
restriction.}
It is tempting to state \eqref{eq:affinity} for the wider class of probes of degree at
most $r-1$, on the reasoning that a probe below degree $r$ cannot contain the degree-$r$
tail character. \emph{That reasoning is wrong}, because the Hellinger distance sees powers
of the probe and degrees add. At $r=3$ take
\[
g(z)=\tfrac{1}{2}\bigl(z_{4}z_{5}+z_{6}\bigr),
\qquad \|g\|_{\infty}=1,\qquad \deg g=2\le r-1 .
\]
With $\chi=z_{4}z_{5}z_{6}$ one computes $g^{2}=\tfrac{1}{2}(1+\chi)$, so
$\mathbb{E}[\chi g^{2}]=\tfrac12$: squaring the probe \emph{recreates} the character the
cap was meant to exclude. On this model that is visible in a bounded statistic,
\[
\mathbb{E}_{+}\bigl[Y^{3}\bigr]-\mathbb{E}_{-}\bigl[Y^{3}\bigr]=3h^{2}
\qquad\text{exactly,}
\]
and since the sixth moments are bounded uniformly for $h\le1$, \eqref{eq:cs} with $f=Y^{3}$
forces $H^{2}\ge c\,h^{4}$ with $c>0$. No uniform $H^{2}\le Ch^{6}$ can hold on the
degree-capped class. Measured on that single probe the order is $3.94$ and
$H^{2}/h^{4}$ settles at $1.74\times10^{-2}$ while $H^{2}/h^{6}$ diverges. \NUM{}

So the class in \eqref{eq:affinity} is the \emph{linear} one, $g(z)=a\cdot z$ with
$\|g\|_{\infty}\le1$, and that is the class Theorem~\ref{thm:vanish} and
\lean{mixture\_hellinger\_rate} are stated for. Figure~\ref{fig:unih} shows the linear
class against the unrestricted one.

\NOTEPAR{A local ascent returns a value \emph{attained} by a legal probe, so it is a lower
bound on the supremum. A lower bound can refute an upper-bound law and can never certify
one. On the degree-capped class it refutes; on the linear class it fails to refute, which
is consistent with the law but is not what establishes it --- the theorem is. An earlier
version of this measurement searched the degree-capped class at random, missed the probe
above, and reported an order of $7.12$; the exhibited probe is the reason the class is now
pinned in the statement rather than chosen for convenience.}

\paragraph{How little nonlinearity destroys the separation.}
The gap between order $2r$ and order $4$ is not a cliff. Widen the class by a single
parameter,
\[
g_{\eta}(z)=\frac{z_{4}+z_{5}+z_{6}+\eta\,z_{4}z_{5}}{3+\eta},
\qquad\|g_{\eta}\|_{\infty}=1,
\]
whose moments against the tail character are exactly
$\mathbb{E}[\chi g_{\eta}]=0$, $\mathbb{E}[\chi g_{\eta}^{2}]=2\eta/(3+\eta)^{2}$ and
$\mathbb{E}[\chi g_{\eta}^{3}]=6/(3+\eta)^{3}$: at $\eta=0$ the signal first appears at the
cube, and any $\eta>0$ opens a square channel. Writing the density difference as
\begin{gather*}
p_{+}-p_{-}=\frac{\eta h^{2}}{(3+\eta)^{2}}A-\frac{h^{3}}{(3+\eta)^{3}}B+R,\\
A=f''(\cdot-1)-f''(\cdot+1),\qquad
B=f'''(\cdot-1)-f'''(\cdot+1),
\end{gather*}
the head density $f$ is even, so $A$ is odd, $B$ is even, the base density is even and
$\int AB/p_{0}=0$: the $\eta h^{5}$ cross term drops out and the two channels are
additive. The measured order in $h$ falls monotonically from $5.94$ at $\eta=0$ through
$5.14$, $4.50$, $4.11$ to $3.95$ at $\eta=1.6$, and the two-channel form
$c_{6}h^{6}+c_{4}\eta^{2}h^{4}$ fitted on the coarse half of the grid predicts the fine
half within its small-$\eta$ regime. Figure~\ref{fig:crossover} shows both. The crossover therefore sits at
$\eta\asymp h$: a nonlinear measurement capability of order $h$ already erases
the advantage. \NUM{}

\NOTEPAR{This is a law for the one probe $g_{\eta}$, and a probe bounds the supremum from
below only. A crossover for a \emph{class} would need a uniform upper bound over that
class, which is not attempted and is not claimed.}

\paragraph{The proof route, now complete.}
The route to \eqref{eq:affinity} that avoids checking Taylor coefficients one at a time has
four steps: separate head from tail; replace the alternating sum over the tail cube by an
$r$-th order finite difference; bound it by $h^{r}\prod|a_{i}|\le h^{r}r^{-r}$; and convert
the density gap into a Hellinger bound against a lower envelope. All four are now proved.

The first step rests on the one-dimensional identity
\[
f(x-b)-f(x+b)=-b\int_{-1}^{1}f'(x-bt)\,dt,
\]
from the fundamental theorem of calculus, and on the exact expansion of the iterate as an
alternating sum over the sign cube,
$D_{b_{1}}\cdots D_{b_{r}}f(x)=2^{-r}\sum_{z\in\{\pm1\}^{r}}\chi(z)f(x-b\cdot z)$.

The $h^{r}$ factor comes from the $r$-fold difference bound, which has to be
\emph{localised}: the global version would need a uniform bound on $\varphi^{(r)}$, which a
Gaussian envelope is not. Only the values of the $r$-th derivative within
$\sum_{i}|b_{i}|$ of the base point matter, and paying for that window widens the centre
range from $R$ to $R+h_{0}$ --- nothing else changes.

The envelope is built from the Hermite-like recursion $\varphi^{(n)}=P_{n}\varphi$ with
$\deg P_{n}\le n$, giving $|\varphi^{(n)}(y)|\le A_{n}(1+|y|)^{n}\varphi(y)$ with $A_{n}$
the coefficient mass. Finiteness of $\int B_{r}^{2}/\ell$ then needs no
polynomial-times-Gaussian integrability lemma at all: $1+x\le e^{x}$ absorbs the
polynomial, completing the square absorbs both linear exponentials --- the one from the
upper envelope and the one from the lower --- and what is left is $e^{-y^{2}/4\sigma^{2}}$.
The conclusion is
\[
\int \frac{\Delta^{2}}{\ell}\;\le\;h^{2r}\,r^{-2r}\,C(\sigma,R+h_{0},r),
\qquad C<\infty \text{ explicit.}
\]
Figure~\ref{fig:envrate} shows both halves: the measured order is $2r$ to two decimals at
$r=2,3,4$, and the envelope ratio sits under its dominating Gaussian at every point of the
grid. \NUM{}

\leanbox{\lean{finite\_difference\_one}, \lean{iterDiff\_eq\_signSum},
\lean{iterDiff\_abs\_le\_window}, \lean{gaussDeriv\_envelope},
\lean{envRatio\_integrable}, \lean{iterDiff\_mixture\_envelope},
\lean{hellinger\_rate\_bound}, \lean{hellinger\_le\_of\_gap} in \lean{LeanPaperBRatesV1}}

\paragraph{The identification the envelope needed, and why its old form was empty.}
The composition above took the density-difference identification as a hypothesis:
$p-q=\mathrm{iterDiff}\,q\,(\text{shifts})$, with $q$ the same mixture on both sides,
together with $p\ge0$. \emph{That pair of hypotheses is contradictory.} With one Gaussian
component and one shift $s>0$ we have $p(y)=q(y)+\tfrac12\bigl(q(y-s)-q(y+s)\bigr)$, and at
$y=-T$ with $T=\frac{s^{2}+4\sigma^{2}}{2s}$ the tail ratio is $q(T-s)/q(T)=e^{2}>3$ while
$q(T+s)\le q(T)$, so $p(-T)<0$. The hypotheses hold only at $h=0$, where the conclusion is
empty: the old composition said nothing wherever it was of interest.

The repair is to stop asking one density to be its own difference, and to build the pair the
difference actually comes from. Write $v(z)$ for the signed shift sum of a sign pattern $z$
and $\mathrm{par}(z)$ for the product of its signs, and set
\begin{gather*}
Q=2^{-r}\sum_{z}M(\cdot-v(z)),\\
P=2^{-r}\sum_{z}\bigl(1+\mathrm{par}(z)\bigr)\,M(\cdot-v(z)).
\end{gather*}
$P$ keeps the even patterns at double weight and drops the odd ones; $Q$ is the flat average.
Both are non-negative combinations of translates of one base, both carry the same total mass,
and $P-Q=\mathrm{iterDiff}\,M\,(\text{shifts})$ \emph{exactly}. The identification is now an
identity about a constructed pair rather than a hypothesis about a hoped-for one, and $p\ge0$
is free. Every centre $c_{j}+v(z)$ lies within $R+h_{0}$, so $Q$ inherits the base's floor and
the envelope bound applies with $\ell=\mathrm{envLo}$:
\[
\int\frac{(P-Q)^{2}}{\mathrm{envLo}}\ \le\ h^{2r}r^{-2r}\,C(\sigma,R+h_{0},r),
\]
with no integrability hypothesis at all. $C$ mentions neither the weights, nor the centres,
nor the shift profile --- one constant, every admissible query, which is the uniformity the
per-query bound owed.

\NOTEPAR{Uniform over the class \emph{as the box constraints define it}: centres within $R$,
$\sum_{i}|ha_{i}|\le h_{0}$, $\sum_{i}|a_{i}|\le1$. That a given observation model's queries
\emph{induce} shifts inside that box is the modelling step and is not proved here. The
constructed pair is unnormalised; the two worlds carry equal mass, so one common normaliser
converts the statement, and the $\chi^{2}$ form needs no integrability hypothesis while the
KL form needs three integrability conditions and $\int(P-Q)=0$, all properties of the pair
rather than modelling assumptions. \emph{Status 2026-09-26: on the parity channel the box is
now discharged from the construction (\lean{queries\_inside\_box}).}}

\leanbox{\lean{self\_difference\_not\_a\_density}, \lean{worlds\_diff},
\lean{worldQ\_floor}, \lean{mixDens\_floor}, \lean{chi2\_uniform},
\lean{kl\_uniform\_over\_queries} in \lean{LeanPaperBWorldsV1}}

\NOTEPAR{Still do not write that the two exponents are proved for continuous adaptive
randomised policies. What the analysis above closes is exactly one link: the per-step
Hellinger bound $H^{2}\le C_{r}h^{2r}$ on the construction, with every constant explicit
and no assumed tail. The link it does \emph{not} close is the one before it. Both
\lean{affinity\_pow\_lower} and \lean{budget\_lower\_from\_affinity} take the transcript
affinity's per-step multiplicative step as a \emph{hypothesis}; supplying that step for an
adaptive randomised policy on a \emph{continuous} observation space is the chain-rule
statement that \lean{LeanPaperBTranscriptKLV1} proves only for a finite alphabet. So the
chain reads: per-step bound (now proved) $\to$ transcript affinity product (proved for
finite alphabets only) $\to$ sample budget (proved arithmetic). The middle link is the
open one, and it is limitation~2).}

A complementary route avoids evaluating any integral: a bounded statistic that separates two
laws in expectation forces a Hellinger separation,
\begin{equation}
\bigl(\mathbb{E}_{P}f-\mathbb{E}_{Q}f\bigr)^{2}\le 4H^{2}(P,Q)\bigl(\mathbb{E}_{P}f^{2}+\mathbb{E}_{Q}f^{2}\bigr),
\label{eq:cs}
\end{equation}
so the exact moment contrasts above give $H^{2}\gtrsim h^{2}$ and $H^{2}\gtrsim h^{2r}$ from
a non-zero leading coefficient alone.

\leanbox{\lean{abs\_pow\_add\_le}, \lean{integrable\_abs\_pow\_gaussian},
\lean{gaussian\_shift\_moment\_le} in \lean{LeanPaperBRatesV1}}

\leanbox{\lean{hellinger\_lower\_of\_moment\_contrast},
\lean{hellinger\_pos\_of\_contrast\_ne\_zero}, with \lean{cauchy\_schwarz}, in
\lean{LeanPaperBHellingerV1}}

\section{The law beyond the witness}
\label{sec:law}

Four worlds are a witness. The general statement is about obstructions, and its
combinatorial half is proved.

\begin{theorem}[Classes resolve iff no obstruction hides]
\label{thm:classes}
For any grouping of worlds into classes, every class has a common safe answer if and only if
no obstruction is contained in a class.
\end{theorem}

\leanbox{\lean{classes\_resolve\_iff}, with \lean{hasCommon\_mono} and
\lean{isObstruction\_mono}, in \lean{LeanUnifiedObstructionOrderV1}}

The classes must be built correctly. We group worlds by equality of their \emph{entire}
coefficient vector up to order $s$, over every admissible query: that is equality of a tuple,
hence an equivalence relation by construction. Taking a pairwise ``indistinguishable at order
$s$'' relation and closing it by union-find is not the same thing and is not used.

\begin{theorem}[Randomised error floor]
\label{thm:floor}
Let $H$ be an obstruction and $p$ any distribution over decisions. Then some $F\in H$
receives unsafe mass at least $1/|H|$.
\end{theorem}

\leanbox{\lean{obstruction\_total\_unsafe\_mass},
\lean{obstruction\_error\_floor\_randomized} in \lean{LeanUnifiedObstructionOrderV1}}

The argument is a fibre count: every decision is unsafe somewhere in $H$, so assigning each
decision to such a world gives disjoint fibres whose masses sum to one. It covers randomised
rules, which is what converts ``the obstruction survived'' into a lower bound on risk.

On the family of Section~\ref{sec:family} at $r=3$, with orders computed from the
$h$-expansion of the real law, the classes behave exactly as the law predicts
(Table~\ref{tab:classes}), giving $r_{\rm safe}=1$ and $r_{\rm id}=3$. \CERT{}

\begin{table}[t]
\caption{Classes by coefficient equality up to order $s$, on the real family at $r=3$.}
\label{tab:classes}
\centering
\footnotesize
\begin{tabular}{@{}rrl@{}}
\toprule
$s$ & classes & every class has a common safe answer \\
\midrule
$0$ & $1$ & no \\
$1$ & $2$ (split by $u$) & yes $\Rightarrow r_{\rm safe}=1$ \\
$3$ & $4$ (split by $v$) & yes $\Rightarrow r_{\rm id}=3$ \\
\bottomrule
\end{tabular}
\end{table}

\TODOPAR{What is \emph{not} proved: that $N_{\rm safe}(h)=\Theta(h^{-2r_{\rm safe}})$ in
general. Section~\ref{sec:law} supplies the combinatorics and the error floor; converting an
order-$s$ indistinguishability into a sample count for an arbitrary family is the analytic
half and is open. The matching orders are established on the family of
Section~\ref{sec:family}, and that is the scope of the rate claim.}

\paragraph{The $\log(1/\delta)$ side, proved.}
A Chebyshev upper bound gives $1/\delta$, not $\log(1/\delta)$, so the two-parameter law
needs its lower side from a change of measure --- and it does \emph{not} need two worlds
with disjoint safe sets. Let $H$ be a class of worlds indistinguishable below the relevant
order with $\bigcap_{\theta\in H}\mathcal{A}_{\varepsilon}(\theta)=\varnothing$, and let
there be $K$ answers. At $\theta$ some answer $a$ is returned with probability at least
$(1-\delta)/K$; because the safe intersection is empty there is $\lambda\in H$ at which $a$
is unsafe, so $\Pr_{\lambda}(\hat a=a)\le\delta$. Data processing onto that single event
and the chain rule then give
\[
N\;\ge\;\frac{\frac{1-\delta}{K}\log(1/\delta)-2e^{-1}}{C\,h^{2s}} .
\]
Every step is proved from scratch: the two-term log-sum inequality from $\log y\ge1-1/y$
alone, hence data processing for a two-cell partition; $x\log x\ge-1/e$; the pigeonhole on
the answers; and the assembly, with the $h$ dependence confined to the denominator so no
constant hides a dependence on either parameter.

\leanbox{\lean{log\_sum\_two}, \lean{mul\_log\_ge\_neg\_inv\_e}, \lean{kl2\_ge\_log},
\lean{exists\_heavy\_answer}, \lean{decision\_budget\_log\_delta},
\lean{decision\_budget\_two\_parameter} in \lean{LeanPaperBRatesV1}}

\paragraph{The headline, assembled --- and which half is a theorem.}
The two exponents bound different quantities, and assembling them is what makes visible
which half is a theorem. On the lower side,
$\mathrm{lowerBudget}$ names the bound above: with a per-step divergence at most
$Ch^{2s}$, any rule correct with probability $1-\delta$ needs
$N\ge\bigl((1-\delta)K^{-1}\log(1/\delta)-2/e\bigr)/C\,h^{-2s}$, the constant depending on
the family and on the number of answers alone --- never on $h$ or $\delta$. It is applied at
$s_{\mathrm{safe}}$ and at $s_{\mathrm{id}}$ with nothing differing but the exponent. On the
upper side, $\mathrm{safeSufficient}$ names the Chebyshev budget, which scales as
$h^{-2}/\alpha$.

The separation between the exponents is then an \emph{identity}, not an estimate: at
$s_{\mathrm{safe}}=1$ and $s_{\mathrm{id}}=t+1$,
\[
\mathrm{lowerBudget}(t+1)
=\frac{36\alpha\,c_{\delta}}{CV}\cdot\frac{1}{h^{2t}}\cdot\mathrm{safeSufficient},
\]
and that factor exceeds any $M$ once $h<\sqrt{A/M}$, with $A$ written down.

\NOTEPAR{The two sides do \emph{not} match in the confidence parameter, and nothing here
pretends otherwise: the lower bound carries $\log(1/\delta)$ and the Chebyshev upper bound
carries $1/\alpha$, with the ratio $VC/36\alpha c_{\alpha}$ recorded as its own statement.
There is no theorem in this corpus of the form $N^{*}=\Theta(h^{-2s}\log(1/\delta))$.
Closing that gap needs a test whose error decays exponentially in $N$ --- the
characteristic-function route --- which is not formalised. \emph{Status 2026-09-26: superseded. The characteristic-function route is now formalized, and
\lean{parity\_rates}, \lean{family\_safe\_minimax} and \lean{family\_id\_minimax} are theorems of
that form. This note concerns the Chebyshev pair only.}}

\leanbox{\lean{budget\_lower}, \lean{safe\_upper}, \lean{separation\_ratio},
\lean{separation\_unbounded}, \lean{sides\_do\_not\_match} in \lean{LeanPaperBHeadlineV1}}

\NOTEPAR{The chain rule from a per-query bound to the transcript is the same open middle
link as in limitation~2): proved for a finite alphabet, not for an adaptive randomised
policy on a continuous observation space. So the display above is a theorem about the
budget \emph{given} that link, and the $\log(1/\delta)$ factor is the part that is new
here.}

\section{Evidence}
\label{sec:evidence}

Every script self-checks against a proved statement and refuses to write output if a check
fails. Evidence classes are printed on every figure and never mixed.

\NOTEPAR{Three claims are in play here and a passing check certifies only one of them at a
time. \textsc{(i) A numerical search} for the probe maximising $H^{2}$: a value attained by
a legal probe is a \emph{lower} bound on the supremum, so it can refute an upper-bound law
and never certify one, and at $r=3$ the search is seed-limited (recorded per $h$).
\textsc{(ii) A theorem uniform over the query class}: the per-step Hellinger bound holds
for every probe in the linear class, which is what \eqref{eq:affinity} consumes --- and
which the search does not establish. \textsc{(iii) A matching rate}: an upper bound and a
lower bound meeting at the same exponent on the same observation model. That is
\emph{not} closed: the upper bound is assembled per pair and the two orders are proved on
\emph{two constructions that are not yet joined}. \texttt{all\_passed} in an artifact
certifies the checks of the script that wrote it. It does not certify (ii) from (i), and it
does not certify (iii) at all. \emph{Status 2026-09-26: (iii) is now closed by the main text's Theorems ``The two rates'',
``Both rates on the latent family'' and ``The general law is tight on both channels''.}}

\begin{table}[t]
\caption{Evidence packages for this paper.}
\label{tab:evidence}
\centering
\footnotesize
\resizebox{\columnwidth}{!}{%
\begin{tabular}{@{}llr@{}}
\toprule
Package & Script & Checks \\
\midrule
Family, safe sets, moment identities & \texttt{EXP\_paperB\_construction} & 8/8 \\
Hellinger orders, real law & \texttt{EXP\_paperB\_hellinger\_real} & 6/6 \\
Per-step bound, over the query class$^{\dagger}$ & \texttt{EXP\_paperB\_uniform\_hellinger} & 6/6 \\
Obstruction law, classes & \texttt{EXP\_paperB\_obstruction\_law} & 9/9 \\
Moment law, degree-graded queries & \texttt{EXP\_paperB\_moment\_tests} & 8/8 \\
Estimator error, Chebyshev budget & \texttt{EXP\_paperB\_estimator\_error} & 7/7 \\
Envelope chain, the (B4) rate & \texttt{EXP\_paperB\_envelope\_rate} & 8/8 \\
Query-class crossover, $g_{\eta}$ & \texttt{EXP\_paperB\_crossover} & 8/8 \\
The two decoders, parity channel & \texttt{EXP\_paperB\_decoder} & 7/7 \\
Query-order vanishing/attainment & \texttt{EXP\_unified\_query\_order} & 6/6 \\
\bottomrule
\end{tabular}}

\vspace{2pt}
{\footnotesize $^{\dagger}$This package measures the per-step bound on the \emph{linear-class}
observation family of Section~\ref{sec:family}. It is \emph{not} evidence about the
constructed pair $(\mathrm{worldP},\mathrm{worldQ})$ of the uniformity theorem, nor about the
decoder's own pair, whose difference is
$D_{b_{1}}\cdots D_{b_{r}}[\varphi(\cdot-1)-\varphi(\cdot+1)]$ by
\lean{channel\_pair\_diff}. An earlier draft cited it as though the three were one family;
they are three, and only the identity between the last two is proved.}
\end{table}

\begin{figure}[t]
\centering
\includegraphics[width=\columnwidth]{paperB_crossover.pdf}
\caption{Left: the measured order in $h$ against $\eta$, falling from six to four. Right:
$H^{2}$ against $h$ at several $\eta$, with the two-channel form
$c_{6}h^{6}+c_{4}\eta^{2}h^{4}$ fitted on the coarse half and drawn through the fine half.
The crossover is at $\eta\asymp h$.}
\label{fig:crossover}
\end{figure}
\begin{figure}[t]
\centering
\includegraphics[width=\columnwidth]{paperB_envelope_rate.pdf}
\caption{The proved envelope chain. Left: $\int\Delta^{2}/\ell$ against $h$ for
$r=2,3,4$, with the measured orders $4.00$, $6.00$, $8.00$ and the proved bound
$h^{2r}r^{-2r}C$ dashed; the vertical gap is the explicit constant, which is not claimed
sharp. Right: the envelope ratio $\mathrm{envUp}^{2}/\mathrm{envLo}$ and the Gaussian
$C\,e^{-y^{2}/4\sigma^{2}}$ that dominates it, which is what makes the integral finite
without any polynomial-times-Gaussian integrability lemma.}
\label{fig:envrate}
\end{figure}
\begin{figure}[t]
\centering
\includegraphics[width=\columnwidth]{paperB_uniform_hellinger.pdf}
\caption{Left: the per-step separation over the query class, for both directions. Right:
the same quantity with every class restriction removed --- the separation collapses, so
the restriction is what creates it. The degree-capped class is not shown: it is refuted in
the text by an exhibited probe, not by a search.}
\label{fig:unih}
\end{figure}

\begin{figure}[t]
\centering
\includegraphics[width=\columnwidth]{paperB_orders.pdf}
\caption{Left: squared Hellinger against $h$. Right: the attained leading coefficient
against $r!/r^{r}$, in exact rational arithmetic.}
\label{fig:orders}
\end{figure}

\begin{figure}[t]
\centering
\includegraphics[width=0.78\columnwidth]{paperB_two_scales.pdf}
\caption{The two scales on one family. Marked as an illustration of the proved orders; the
empirical sample counts are reported separately.}
\label{fig:scales}
\end{figure}

\section{The two orders, on the real geometry, and a matching upper bound}
\label{sec:geometry}

The separation this paper is about is between two \emph{observation orders}: the order at
which a decision becomes safe, and the order at which the world becomes identified. Both
are defined on the real \lean{selectedRadius} geometry rather than on a stand-in.

Write $\theta\sim_{s}\lambda$ for indistinguishability at order $s$, and
$\mathrm{class}_{s}(\theta)$ for the class it induces. Two facts pin the orders.

\begin{theorem}[The two orders separate]
\label{thm:orders}
For $m\ge3$: every class at order $1$ has a common safe answer, no class at order $0$ does,
and identification holds at order $s$ if and only if $s\ge m-1$. Hence
$s_{\mathrm{safe}}=1$ while $s_{\mathrm{id}}=m-1$, and the gap grows with $m$.
\end{theorem}

The identification half is an \emph{iff}, so $m-1$ is tight on both sides rather than an
upper estimate. Two further facts say the geometry is the interesting one and not a
degenerate case: any two worlds share a safe answer at $m\ge3$
(\lean{safe\_sets\_overlap}), while no answer is safe for all of them
(\lean{no\_common\_safe\_all}). Overlapping-but-not-nested safe sets are exactly what makes
the safe order small and the identification order large. The order-$0$ class is an
obstruction (\lean{obstruction\_at\_order\_zero}), which is the hypothesis the
$\log(1/\delta)$ lower bound consumes, so the lower side is applied to a geometry that
verifiably satisfies it.

\leanbox{\lean{identifiableAt\_iff}, \lean{decidableAt\_of\_one\_le},
\lean{not\_decidableAt\_zero}, \lean{decide\_order\_lt\_identify\_order},
\lean{safe\_sets\_overlap}, \lean{no\_common\_safe\_all},
\lean{obstruction\_at\_order\_zero} in \lean{LeanPaperBGeometryV1}}

\paragraph{The partition is the observational one, on an explicit moment family.}
$\sim_{s}$ was \emph{defined} on the world index --- ``equal, or both at index at least
$s$'' --- so Theorem~\ref{thm:orders} was conditional on that being the relation an
observer faces. On the family whose world $i$ has its first non-vanishing moment at order
$i+1$, it is the observational relation, in both directions. Call $q$ a query of order at
most $s$ when it is a linear functional of the moments of order $1,\dots,s$. Then

\begin{theorem}[The order-$s$ classes are exactly the query-blind pairs]
\label{thm:indist}
For every $s$ and every pair of worlds $\theta,\lambda$,
\begin{gather*}
\theta\sim_{s}\lambda
\quad\Longleftrightarrow\\
q(\theta)=q(\lambda)\ \text{ for every query } q \text{ of order at most } s .
\end{gather*}
\end{theorem}

Left to right is termwise: such a $q$ is supported where the two moment vectors agree.
Right to left is a construction: if $\theta\not\sim_{s}\lambda$ then $\theta\ne\lambda$
and one of them, say $\theta$, has index below $s$, and the indicator at order $\theta+1$
is a query of order at most $s$ taking the value $1$ at $\theta$ and $0$ at $\lambda$.

\leanbox{\lean{agreeUpTo\_iff\_indistAt}, \lean{query\_blind\_of\_indistAt},
\lean{exists\_query\_separating}, \lean{indistAt\_iff\_no\_separating\_query} in
\lean{LeanPaperBIndistV1}}

\paragraph{The table, realised by densities.}
Read as a table of raw moments the coefficient table is not realisable by any law --- world
$0$ would need $\mathbb EX=1$ and $\mathbb EX^{2}=0$. It is meant as the table of expansion
coefficients in the perturbation scale, and until a family of densities realises it the
partition rests on a table nobody has exhibited. So we exhibit one. For any density
$\varphi$, put
\[
p_{i}(h)\;=\;(1-h^{i+1})\,\varphi\;+\;h^{i+1}\,\varphi(\cdot-1).
\]
Each $p_{i}(h)$ is a genuine probability density for $0\le h\le1$ --- a two-point mixture, so
non-negativity and unit mass need no delicate integral --- and its expansion has coefficient
$[k=i+1]\cdot(\varphi(\cdot-1)-\varphi)$ at order $k\ge1$, which is the table with one fixed
profile.

The resolving order then follows \emph{from the densities}, in both directions:
$|p_{i}(h)-p_{j}(h)|\le h^{\min(i,j)+1}|\text{profile}|$ everywhere, so nothing separates
earlier; and $\ge\tfrac12h^{\min(i,j)+1}|\text{profile}|$ at a point once $h$ is small, so
something separates there. Hence $r(i,j)=\min(i,j)+1$, and at a family whose only unsafe
subset is the whole world set this gives $s_{\mathrm{safe}}=1$ and $s_{\mathrm{id}}=m-1$.

\NOTEPAR{This settles that the reading is consistent and supplies the orders both ways from
actual densities. It does \emph{not} derive the parity-mask construction's own observation
law; that its densities carry exactly these coefficients is the remaining modelling step, and
until it is written the two exponents are established for the exhibited family and not for
that construction. \emph{Status 2026-09-26: superseded. The parity channel's own graded $\chi^2$ bound is proved
(\lean{parChan\_chi\_le}, \lean{parity\_graded}).}}

\leanbox{\lean{dens\_integral\_one}, \lean{dens\_vanishes\_below},
\lean{dens\_separates\_at}, \lean{coefOf\_realised}, \lean{s\_safe\_eq\_one},
\lean{s\_id\_eq} in \lean{LeanPaperBDensityV1}}

\NOTEPAR{What Theorem~\ref{thm:orders} is and is not, after Theorem~\ref{thm:indist}. The
partition is no longer posited, but two things stand between
$s_{\mathrm{safe}}=1$, $s_{\mathrm{id}}=m-1$ and a pair of \emph{sample} complexities.
First, the moment vectors of Theorem~\ref{thm:indist} are \emph{prescribed}: no
distribution realising them is constructed, and joining this to the moment-matching laws of
Section~\ref{sec:law} is what would make it one family with $F_{\theta}$, a support, a
loss and an observation law. Second, the safe sets here are the real \lean{selectedRadius}
ones, but the obstruction experiment's four $(u,v)$ worlds have safe sets that are
\emph{equal} within $u$ and disjoint across it --- a different family from the intersecting
one this theorem is about, and the two must not be joined in prose. The two orders should
still not be quoted as sample complexities. \emph{Status 2026-09-26: superseded on the parity channel and the latent family, where the
orders are now sample complexities (main text).}}

\subsection{The safe-observation-order law}
\label{sec:safelaw}

The question the paper is for: how deep must an observer resolve worlds before every class
it can still confuse admits one answer that is safe across it? Write $\mathrm{sep}(\theta,
\lambda)$ for the first order at which some admissible query separates two worlds. Call a
set of worlds \emph{together at level $s$} when no pair in it is resolved by order $s$, and
\emph{unsafe} when its worlds share no safe answer.

\begin{theorem}[The safe order is a statement about obstructions]
\label{thm:safelaw}
Order $s$ suffices --- every set together at level $s$ has a common safe answer --- if and
only if every unsafe set already contains a pair resolved by order $s$. Equivalently
\[
s_{\rm safe}=\max_{H\ \text{unsafe}}\ \min_{\theta\ne\lambda\in H}\mathrm{sep}(\theta,\lambda).
\]
\end{theorem}

Quantifying over \emph{all} unsafe sets rather than minimal ones is not a weakening:
unsafety is upward closed and enlarging a set can only lower its pairwise minimum, so the
maximum is attained on minimal obstructions anyway --- and stating it this way removes the
need to define minimality. Two edge cases the law names: if no set is unsafe then order zero
suffices, the decision costing no observation at all; and the property is monotone in $s$.

\paragraph{The overlapping construction, geometry half.}
With at least three worlds and $\mathrm{safe}(w)$ every answer but $w$: each pair of safe
sets meets, the total intersection is empty, and the whole world set is the \emph{only}
unsafe set. Theorem~\ref{thm:safelaw} then reads
$\mathrm{SafeAt}(s)\iff\min_{\theta,\lambda}\mathrm{sep}(\theta,\lambda)\le s$, the
simplest instance it has.

\NOTEPAR{What is an input here, and it is the substance. $\mathrm{sep}$ is \emph{given}.
Deriving it from the densities of an actual observation model --- the coefficient of $h^{k}$
in the expansion, uniformly over the admissible query class --- is this paper's central
obligation and is not done. The law says what follows \emph{from} a resolving-order table,
not that the model has one; the parity-mask query showing that query $j$ resolves world $j$
at order $j$ is the construction's other half and is likewise unproved. And at this family a
query-free guess already has worst-case error $1/m$, so any positive sample-complexity lower
bound must restrict $\delta$ below that floor --- we do not state one that does not. \emph{Status 2026-09-26: superseded. On the parity channel $\mathrm{sep}$ is read off the
channel's mask (\lean{channel\_sep\_is\_read\_off}); query $j$ separates world $j$ at order $j$
through the characteristic-function gap (\lean{query\_separates\_iff}, \lean{parity\_gap\_order});
and the order is established as a graded $\chi^2$ upper bound (\lean{parChan\_chi\_le}) with a
Hellinger lower bound at the separating query (\lean{parity\_hel\_sep}). An exact coefficient of
$h^{k}$ for arbitrary families is not claimed.}}

\leanbox{\lean{safeAt\_iff}, \lean{unsafe\_mono}, \lean{safeAt\_mono},
\lean{safeAt\_zero\_of\_no\_unsafe}, \lean{coSafe\_pairwise\_meet},
\lean{coSafe\_univ\_unsafe}, \lean{coSafe\_proper\_safe}, \lean{coSafe\_safeAt\_iff} in
\lean{LeanPaperBSafeOrderV1}}

\subsection{From an envelope to a divergence bound}
\label{sec:envelope2kl}

The lower bound needs $\mathrm{KL}$, and the envelope work bounds
$\int(p-q)^{2}/\ell$. A Hellinger bound does not give a $\mathrm{KL}$ bound, and neither
does the envelope by itself; one extra inequality is needed, and it is elementary.

\begin{theorem}[Envelope to divergence]
\label{thm:env2kl}
For $p\ge0$ and $q>0$,
$p\log(p/q)\le(p-q)^{2}/q+(p-q)$, the right-hand side being an \emph{identity} for
$p(p/q)-p$; and $(p-q)^{2}/q\le(p-q)^{2}/\ell$ for any floor $0<\ell\le q$. Integrating
against a common measure, the linear term contributes $\int p-\int q=0$ for two densities,
so
\[
\mathrm{KL}(p\|q)\;\le\;\int\frac{(p-q)^{2}}{\ell}.
\]
\end{theorem}

\NOTEPAR{Why the Hellinger route does not substitute, stated as arithmetic rather than as a
remark: the chi-squared integrand exceeds the Hellinger one by $(\sqrt p+\sqrt q)^{2}/q$,
which is unbounded as $p/q$ grows. What remains, and it is the model-specific work: the
density-difference identity, an envelope $|p-q|\le h^{s}B_{s}(y)$ \emph{uniform in the
query}, and a floor with $\int B_{s}^{2}/\ell<\infty$ --- together with the integrability
hypotheses the integral form above carries explicitly. \emph{Status 2026-09-26: done for the constructed pair and the latent family
(\lean{worlds\_diff}, \lean{chi2\_uniform}, \lean{famChi\_v\_le}, \lean{famChi\_u\_le}).}}

\leanbox{\lean{kl\_integrand\_le}, \lean{chi\_integrand\_mono},
\lean{kl\_integrand\_le\_floor}, \lean{kl\_le\_envelope},
\lean{hellinger\_integrand\_le\_chi} in \lean{LeanPaperBKLV1}}

\subsection{The decoder, and the uniformity a class test needs}
\label{sec:decoder}

An upper bound needs a decoder, not a collection of pairwise error bounds. Two things have to
be right, and the second is easy to skip.

A test between two \emph{classes} is not a test between two representatives: what must hold
is $\sup_{\theta\in C}\Pr_{\theta}(\text{says }D)\le\delta'$ and symmetrically,
\emph{uniformly over each class}. Choosing one world from each and testing those two says
nothing about the rest.

\begin{theorem}[Uniform composite test]
\label{thm:composite}
Let $m(\theta)$ be a statistic's mean at $\theta$. If $|m(\theta)-\mu_{C}|\le R$ for
\emph{every} $\theta\in C$, $|m(\lambda)-\mu_{D}|\le R$ for every $\lambda\in D$, and
$\mu_{C}+2R+2\,\mathrm{tol}<\mu_{D}$, then the threshold at $(\mu_{C}+\mu_{D})/2$ places
every world of $C$ strictly below it and every world of $D$ strictly above, whenever the
empirical mean is within $\mathrm{tol}$ of its own world's mean.
\end{theorem}

The only probabilistic input is per-world concentration, through $\mathrm{tol}$; the class
uniformity is carried by $R$, not by a representative.

The output must also exist. After elimination the survivors are still together at level $s$,
and Theorem~\ref{thm:safelaw} is exactly the statement that such a set shares a safe answer
--- so the decoder has something to return, and since the truth is among the survivors, what
it returns is safe at the truth. That is the step a pairwise-error accounting does not give:
correctness of the \emph{output}, not of the tests.

\NOTEPAR{What this does not include. The concentration itself (a bounded statistic and
Hoeffding), the queries realising the separation at each order, and the \emph{total} sample
budget. The number of distinct queries, the number of tests, and the per-test sample count
are three different quantities; we do not report one in place of another, and none of the
three is established here. \emph{Status 2026-09-26: superseded by \lean{safe\_cost}, \lean{id\_cost} in
\lean{LeanPaperBParityDecoderV1} and \lean{safe\_decoder\_error}, \lean{id\_decoder\_error} in
\lean{LeanPaperBFamilyDecoderV1}.}}

\leanbox{\lean{composite\_test\_low}, \lean{composite\_test\_high},
\lean{composite\_test\_separates}, \lean{separation\_of\_gap},
\lean{decoder\_output\_exists}, \lean{decoder\_is\_safe\_at\_truth} in
\lean{LeanPaperBDecoderV1}}

\paragraph{The lower bound at an unresolved class --- and the floor it is conditional on.}
Take a class $H$ still unresolved at order $s-1$ whose worlds share no safe answer. Under
$\theta_{0}\in H$ some answer is returned with probability at least $(1-\delta)/K$; it is
wrong at some $\lambda\in H$; and data processing on the binary event
$\{\widehat a = a\}$ turns the transcript divergence into a bound on the binary divergence,
which the budget arithmetic converts into $\Omega(h^{-2s}\log(1/\delta))$.

The bound is conditional on $\delta$, and not as a technicality. With
$\mathcal A_{w}=[m]\setminus\{w\}$ a decoder that asks \emph{nothing} and guesses uniformly is
unsafe with probability exactly $1/m$. So a claim of the form ``$N=\Omega(h^{-2s}\log(1/\delta))$
for every $\delta<1/2$'' is \emph{false} for this geometry, and the theorem restricts $\delta$
below the floor.

On the decoder side, elimination needs two inclusions --- the truth survives, and nothing
outside its class does --- and the output then exists by the definition of
$s_{\mathrm{safe}}$. The cost is the \emph{sum} over tests, not the per-test sample size.

\leanbox{\lean{decision\_lower\_from\_class}, \lean{zero\_query\_floor},
\lean{no\_lower\_bound\_above\_floor}, \lean{candidate\_contains\_truth},
\lean{candidate\_within\_class}, \lean{decoder\_output\_exists\_of\_common},
\lean{union\_bound\_tests} in \lean{LeanPaperBDecideV1}}

\paragraph{Attainment by a bounded statistic.}
On $(0,1]$ we have $\sin x\ge 3x/4$, so the tail factor's $r$-th power is at least
$(3x/4)^{r}$: the mean gap is $\Theta(h^{r})$ with an explicit constant, a theorem about the
characteristic function rather than a fitted slope. The statistic $\cos(t_{0}Y)$ or
$\sin(t_{0}Y)$ is bounded by one whatever $r$ is, so its variance is at most one --- which is
the whole reason to use it, since the $Y^{r+1}$ moment estimator's variance grows with $r$.
The midpoint threshold then classifies correctly once both empirical means are within half the
gap.

\NOTEPAR{An earlier version of this paragraph said a Hoeffding-type inequality was not
available in the library we formalise against, and stopped at Chebyshev's $1/\delta$. That was
wrong: Hoeffding's lemma and the Hoeffding inequality for independent sub-Gaussian summands
are both present. Since $|f|\le 1$, the centred statistic is sub-Gaussian with proxy $1$, the
midpoint test errs with probability at most $e^{-Ng^{2}/8}$, and $N\ge 8\log(1/\delta)/g^{2}$
suffices; with $J$ tests, $N\ge 8\log(J/\delta)/g^{2}$. The ratio of the two sides of the
headline is then bounded by a constant free of $\delta$. What still limits the claim is not
the confidence parameter but the evidence: three $(m,h)$ cells, all at $m=3$, certify both
decoders, so there is no measured identification-rate curve.}

\leanbox{\lean{sin\_lower\_small}, \lean{gap\_at\_scale},
\lean{bounded\_statistic\_variance}, \lean{threshold\_decoder\_correct},
\lean{budget\_from\_gap} in \lean{LeanPaperBAttainV1}}

\paragraph{The parity-mask channel, and the orders read off it.}
$K_{w}=\{0,\mathrm{tip}_{w}\}$ gives $\mathrm{tip}_{j}\in K_{w}$ exactly when $j=w$, so the
mask $v(j,w)=1-2\mathbf1\{j=w\}$ is $-1$ precisely at the named world. Query $j$ returns
\[
Y_{j}=v(j,w)\prod_{\ell\le j}Z_{\ell}+\frac hj\sum_{\ell\le j}Z_{\ell}+\xi ,
\]
a Gaussian mixture over the sign cube. The cube identity \lean{parity\_cf\_difference} gives
the \emph{unnormalised} difference over the $2^{j}$ cube points,
$2^{j+1}\sinh(it)\sinh(ith/j)^{j}$; it is not yet a characteristic function. The cube is
uniform, so each point carries mass $2^{-j}$, and the additive noise contributes its own
Gaussian factor. Supplying both, the characteristic-function difference of the two
\emph{observation} laws is
\[
\varphi_{+,j}(t)-\varphi_{-,j}(t)
=2\,i^{\,j+1}\sin t\,\bigl(\sin (th/j)\bigr)^{j}e^{-\sigma^{2}t^{2}/2},
\]
of magnitude $\Theta(h^{j})$ --- the leading constant is $2$, not $2^{j+1}$, because the
$2^{-j}$ normalisation cancels all but one power of two. The decoder experiment implements
this normalised form, so no measured number depends on the distinction; the $h$-dependence,
which is what the order claim uses, is the $j$-th power either way. The $v$-dependent term is real for odd $j$ and imaginary for even
$j$, which is why the statistic is chosen by parity. Query $j$ therefore decides $w=j$ at
order $j$, a pair resolves at $\min(w,w')$, and the classes at level $s$ are
$\{1\},\dots,\{s\},\{s+1,\dots,m\}$: $s_{\mathrm{safe}}=1$ and $s_{\mathrm{id}}=m-1$.

\NOTEPAR{This is a \emph{second} observation channel, on support membership. It is not the
raw-$F$ query law of the four-world family, and both are instantiations of the same order law;
writing it as one model would be wrong.}

\leanbox{\lean{maskVal\_eq\_neg\_one\_iff}, \lean{parity\_cf\_difference},
\lean{parity\_gap\_order}, \lean{query\_separates\_iff}, \lean{parity\_s\_safe},
\lean{parity\_s\_id}, \lean{parity\_classes} in \lean{LeanPaperBParityV1}}

\begin{figure*}[t]
\centering
\includegraphics[width=\textwidth]{paperB_decoder.pdf}
\caption{The two decoders on one channel. Left and centre: worst-world error against the
\emph{total} budget --- for identification that is the sum over the $m-1$ queries, not the
per-test sample size --- with Clopper--Pearson bounds from the saved counts (\texttt{paperB\_decoder\_grid\_v2}; no
resampling). A budget counts as achieving $\delta$ only when the upper bound clears it. The
reported budget of a panel is the first rung of an eleven-rung ladder that clears the target,
which is a selection over $J=616$ intervals, so BOTH tails of every interval are budgeted at
$0.05/(2J)$ and all bounds hold simultaneously at $95\%$. Eleven of the sixteen panels have an
upper certificate, six of those also have a failed prefix, and five have none. THREE $(m,h)$
cells certify both decoders, all at $m=3$: at $h=0.6$, $0.45$, $0.32$ the total budgets are
$25\to6400$, $100\to12800$ and $100\to51200$ --- ratios $256$, $128$, $512$; five of the
eight identification panels remain censored. So the panel shows that the safe decoder is
cheaper where both certify, and is \emph{not} an identification-rate curve.
A certified budget is an \emph{upper} bound and nothing more: a rung that failed to certify
is a rung whose upper bound came out above $\delta$, which is \emph{not} evidence that its
true error exceeds $\delta$. Nor is a fixed decoder's error known to be monotone in $N$, so
one failing rung below a success excludes nothing between them --- only a \emph{contiguous
prefix} of tested budgets whose LOWER bounds all exceed $\delta$ bounds the first successful
\emph{tested} budget from below (dotted curves). That bracket concerns this grid; it bounds
neither the untested integers between rungs nor the minimum over all decoders. At $m=3$ the
identification slope interval is $[-5.59,\,-1.07]$, which contains both $-2(m-1)=-4$ and the
safe decoder's $-2$: this panel establishes an achieved-cost separation and \emph{cannot}
separate the two exponents. The ratios above are ratios of rungs meeting the ladder's
standard, not measured complexity ratios. Right: the log-log slope of the mean gap
converges to the query index $j$ as $h$ shrinks; the coarse grid sits below it because of the
sine's cubic term, and the drag is shown rather than hidden.}
\label{fig:bdecoder}
\end{figure*}

\NOTEPAR{Which identification decoder the counts measure (corrected 2026-09-24). The script's
identification rule returns no answer, counted as an error, when two or more tests fire; the
formalized decoder returns a default world there. On every outcome vector and at every world a
wrong formalized guess implies a wrong script guess, so the script's counts are an UPPER bound
on the formalized decoder's error, and its upper certificates are conservative for it. Its
LOWER confidence bounds are not: a lower bound on the script rule's error does not bound the
formalized decoder's error from below. An earlier version of this note called the script's
lower brackets ``the conservative side''; that was wrong. Both bounds are now computed from the
formalized decoder's own counts (the paired rerun \texttt{EXP\_paperB\_decoder\_leanfallback},
same seed and draws, 10/10), and they coincide with the script rule's on all $16$ panels
(\texttt{FIG\_paperB\_decoder\_formal}, 4/4). The figure of the main text is drawn from those
counts.}

\leanbox{\lean{pyGuess}, \lean{decodeB\_of\_pyGuess}, \lean{lean\_wrong\_imp\_py\_wrong} in
\lean{LeanPaperBDecoderMatchV1}}

\subsection{The upper bound, with its per-pair input derived}

An upper bound that \emph{assumes} a per-pair test bound assumes the result. Here the test
is exhibited and its error computed. Let $P,Q$ be the $N$-fold product laws of two
candidate worlds on a finite observation alphabet, and let the test declare the first world
exactly on the sequences it makes at least as likely. Its two error probabilities satisfy
\begin{equation}
\tfrac12\bigl(\mathrm{err}_{\mathrm{I}}+\mathrm{err}_{\mathrm{II}}\bigr)
=\tfrac12\sum_{t}\min\bigl(P(t),Q(t)\bigr),
\label{eq:lrtexact}
\end{equation}
an \emph{identity}, not a bound: the two sums recombine term by term. Bounding the minimum
by the geometric mean and tensorising the Hellinger affinity,
$\sum_{t}\sqrt{P(t)Q(t)}=\bigl(\sum_{o}\sqrt{p(o)q(o)}\bigr)^{N}=(1-H^{2})^{N}$, gives

\begin{theorem}[Per-pair test bound]
\label{thm:lrt}
If $H^{2}(p,q)\ge c$ then \emph{each} of the test's two error probabilities is at most
$\delta/K$ once $N\ge\log(K/\delta)/c$; over $K$ pairs each sums to at most $\delta$.
\end{theorem}

\begin{proof}
By \eqref{eq:lrtexact} the two errors sum to $\sum_{t}\min(P(t),Q(t))$. Since
$\min(a,b)^{2}\le ab$ for $a,b\ge0$, each summand is at most $\sqrt{P(t)Q(t)}$. Writing
$P(t)=\prod_{i}p(t_{i})$ and likewise for $Q$, the square root of the product is the
product of the square roots, and summing a product over all sequences factorises:
\begin{align*}
\sum_{t}\sqrt{P(t)Q(t)}
&=\sum_{t}\prod_{i}\sqrt{p(t_{i})q(t_{i})}\\
&=\Bigl(\sum_{o}\sqrt{p(o)q(o)}\Bigr)^{N}
=(1-H^{2})^{N}.
\end{align*}
Now $1-x\le e^{-x}$ gives $(1-H^{2})^{N}\le e^{-NH^{2}}\le e^{-Nc}$, and
$N\ge\log(K/\delta)/c$ makes that at most $\delta/K$. Both error probabilities are
non-negative and their sum is bounded, so each is bounded separately --- which is the form
the union bound over the $K$ pairs consumes. Note the bound is on the sum, not on the
average: an earlier version halved it, which bounds neither world's own error.
\end{proof}

With $c=c_{0}h^{2r}$ this is $N=O\bigl(h^{-2r}\log(K/\delta)\bigr)$: the \emph{same}
exponent in $h$ that the lower bound produces, so the two sides match in order. The
$\delta$-dependence matches only up to a constant, which \lean{bounds\_same\_order} makes
explicit. The separation $c_{0}h^{2r}\le H^{2}$ is an input to this theorem --- it is what
Section~\ref{sec:law} supplies on the linear class --- and here the two are not composed into
a single statement (the composition is \lean{hellinger\_tournament} and
\lean{share\_tournament}); \lean{separation\_nonvacuous} exhibits laws with $0<H^{2}<1$ so
the theorem is not vacuous.

\leanbox{\lean{affinity\_prodLaw}, \lean{err\_eq\_sum\_min},
\lean{err\_le\_affinity\_pow}, \lean{lrt\_error\_le},
\lean{identify\_sample\_complexity}, \lean{decoder\_total\_error},
\lean{separation\_nonvacuous} in \lean{LeanPaperBTestV1}; \lean{decoder\_union\_bound},
\lean{bounds\_same\_order} in \lean{LeanPaperBGeometryV1}}

\subsection{The endpoint on the parity channel: both rates, every almost-surely stopping policy}
\label{sec:endpoint}

On the parity-mask channel the pieces above are composed into one statement about the
minimax cost itself. It uses the continuous Gaussian observation model and adaptive
randomised policies, so it does not rest on the finite-alphabet transcript or the
finite-alphabet test. Write $N^{*}_{\mathrm{safe}}(h,\delta)$ and $N^{*}_{\mathrm{id}}(h,\delta)$
for the minimax worst-world expected number of non-dummy queries. The minimum is over
$\delta$-correct policies: every world has non-termination probability zero and error
probability at most $\delta$ for the task.

\begin{theorem}[The two rates on one channel]
\label{thm:endpoint}
Fix $m\ge3$ and noise $\sigma>0$. Take the parity channel with the tip loss and tolerance
$0$, so a safe answer at world $w$ is any $a\ne w$. The policy class is every randomised
adaptive policy with a measurable stopping rule that stops almost surely, and with
measurable answer and guess read at the stopping time. The stopping time need not be
bounded, and its expectation need not be finite. There are
$0<c\le C$, $h_{0}>0$ and $\delta_{0}\in(0,1)$, depending on $m$ and $\sigma$ only, such that
for all $0<h\le h_{0}$ and $0<\delta\le\delta_{0}$
\begin{gather*}
c\,\frac{\log(1/\delta)}{h^{2}}\le N^{*}_{\mathrm{safe}}(h,\delta)\le C\,\frac{\log(1/\delta)}{h^{2}},\\
c\,\frac{\log(1/\delta)}{h^{2(m-1)}}\le N^{*}_{\mathrm{id}}(h,\delta)\le C\,\frac{\log(1/\delta)}{h^{2(m-1)}} .
\end{gather*}
\end{theorem}

The theorem is proved first for policies with a finite horizon, and then carried to the
whole class by truncation. Run the stopped policy, which issues the free dummy query after
stopping, for $n$ rounds, with $n$ chosen so that every world has stopped with probability at
least $1-\delta$. The result is a finite-horizon policy that is correct at level $2\delta$ and
costs no more than the original. The finite-horizon lower bound applied to it gives the same
order, at level $2\delta$. Conversely, a finite-horizon decoder is a stopping policy that
stops at its horizon, with the same cost. So the upper side, and the non-vacuity of the
class, transfer too. Both lower sides hold against every policy in the class. A correct safe policy's answer at
one world puts mass $(1-\delta)/(m-1)$ on an event of mass at most $\delta$ at another world,
which forces the divergence to be at least $\log(1/\delta)/4(m-1)$. The transcript's
divergence is at most the per-query $\chi^{2}$, of order $h^{2}$, times the expected number of
queries. For identification the pair of worlds is named only by query $m-1$, whose
information is of order $h^{2(m-1)}$. The upper sides are the two decoders of
Figure~\ref{fig:bdecoder}. They use the bounded statistic $\cos(t_{0}Y)$ or $\sin(t_{0}Y)$,
per-test budgets from a concentration bound, and the total budget summed over every test the
decoder runs.

\leanbox{\lean{parity\_tip\_finish}, \lean{parity\_rates}, \lean{safeCert},
\lean{idCert} in \lean{LeanPaperBParityFinishV1}; \lean{safeLower}, \lean{idLower} in
\lean{LeanPaperBParityExperimentV1}; \lean{transcript\_kl\_le} in
\lean{LeanPaperBProtocolV1}; \lean{parChan\_chi\_le} in \lean{LeanPaperBParityChannelV1};
\lean{event\_klDiv\_ge}, \lean{safe\_pair\_exists} in \lean{LeanPaperBObstructionV1};
\lean{tip\_finish\_from\_protocols} in \lean{LeanPaperBRateClosureV1};
\lean{stopped\_parity\_rates}, \lean{stopped\_lower\_log}, \lean{stopped\_safe\_lower},
\lean{stopped\_id\_lower},
\lean{trunc\_correct}, \lean{outT\_first}, \lean{ofPol\_correctS}, \lean{costS\_ofPol} in
\lean{LeanPaperBStoppedParityV1}}

\NOTEPAR{Scope. The cost of a stopping policy is the supremum over horizons of the expected
number of non-dummy queries, which is $\mathbb{E}\tau$ by monotone convergence. Almost-sure
stopping is required uniformly over the $m$ worlds: for every $\varepsilon>0$ one horizon
at which every world has stopped with probability at least $1-\varepsilon$. The lower bounds
over this class are at level $2\delta$ and hold for $2\delta\le\min(\delta_{S},\delta_{I})$,
which changes the constants and not the orders. A policy stops by issuing the
zero-information dummy query, which is not counted. The constants are uniform in $h$ and $\delta$ on the rectangle,
but depend on $m$ and $\sigma$. The theorem concerns this channel and this loss, not the
general law of Section~\ref{sec:law}. The decoder budgets are the formalized decoders'
own, established by a paired rerun (see the note after Figure~\ref{fig:bdecoder}).}

\subsection{The crossover: the mechanism, not a sweep}
\label{sec:crossover}

Widening the probe to $g_{\eta}=(z_{4}+z_{5}+z_{6}+\eta z_{4}z_{5})/(3+\eta)$ is expected
to give $H^{2}\asymp h^{6}+\eta^{2}h^{4}$. A sweep of fitted slopes cannot decide that
shape, because it cannot see whether the two scales \emph{add} or \emph{interfere}: that
turns on a cross term. The algebra settles it.

Against the hidden character $\chi=z_{4}z_{5}z_{6}$,
\begin{gather*}
\mathbb{E}[\chi g_{\eta}]=0,\qquad
\mathbb{E}[\chi g_{\eta}^{2}]=\frac{2\eta}{(3+\eta)^{2}},\\
\mathbb{E}[\chi g_{\eta}^{3}]=\frac{6}{(3+\eta)^{3}},
\end{gather*}
all three exact. The density difference between $v=+1$ and $v=-1$ then expands with an odd
channel $A(y)=f''(y-1)-f''(y+1)$ at coefficient $\eta h^{2}/(3+\eta)^{2}$ and an even
channel $B(y)=f'''(y-1)-f'''(y+1)$ at coefficient $h^{3}/(3+\eta)^{3}$.

\begin{theorem}[The cross term vanishes by parity]
\label{thm:parity}
If $A$ is odd, $B$ is even and the base density $p_{0}$ is even, then $AB/p_{0}$ is odd and
$\int AB/p_{0}=0$. Consequently
$\int(\alpha A+\beta B)^{2}/p_{0}=\alpha^{2}\int A^{2}/p_{0}+\beta^{2}\int B^{2}/p_{0}$,
and substituting the two coefficients gives exactly $\eta^{2}h^{4}$ plus $h^{6}$.
\end{theorem}

Numerically, on the model's own baseline --- the mixture $p_{0}(y)=\tfrac12[f(y-1)+f(y+1)]$
over the $z_{1}$ channel, not the head density $f$ alone, which an earlier version of the
check used --- the normalised cross term is $6.6\times10^{-15}$.

\NOTEPAR{Three things stand between this mechanism and the law, and none is proved: that
$A$ and $B$ are not identically zero, that the two integrals are finite, and that the
remainder is bounded \emph{uniformly in} $(h,\eta)$. It is also a law for the single probe
$g_{\eta}$; a class-level crossover needs the widened class defined and a uniform upper
bound over it. The measured orders $5.93\to3.95$ are evidence, not proof.}

\leanbox{\lean{moment\_one}, \lean{moment\_two}, \lean{moment\_three},
\lean{integral\_odd\_eq\_zero}, \lean{cross\_term\_vanishes},
\lean{quadratic\_splits}, \lean{two\_scale\_shape}, \lean{crossover\_balance} in
\lean{LeanPaperBCrossoverV1}}


\section{The general law's newer sides, and the latent family}
\label{sec:newer}

\paragraph{Unbounded stopping (Theorem ``The lower law for unbounded stopping'')} A stopping
policy has query kernels for every round, a measurable stop set after every number of rounds, and
an output read at the stopping time. After its stop it issues the dummy, which is free and
carries no information. One kernel policy therefore serves every horizon, and the law of its
first $n$ rounds is the $n$-round history law. The output at horizon $n$ is the output at the
first stopping time, if the stop has fired. A horizon at which every world has stopped with
probability at least $1-\delta$ makes the truncation correct at level $2\delta$. Its cost is at most
the supremum of the truncated costs. A finite-horizon policy embeds as a stopping policy with the
same cost, since after its horizon it spends nothing.

\leanbox{\lean{hRestr}, \lean{fired}, \lean{stoppedPol}, \lean{outT}, \lean{outT\_first},
\lean{hist\_congr}, \lean{hist\_map\_restr}, \lean{costSt}, \lean{CorrectSt}, \lean{trunc\_correct},
\lean{trunc\_two}, \lean{stopped\_obstruction\_lower\_chi}, \lean{ofFinite},
\lean{costSt\_ofFinite}, \lean{correctSt\_ofFinite} in \lean{LeanPaperBStoppedLawV1}}

\paragraph{The Hellinger tournament (Theorem ``The upper law on every channel'')} The
Bhattacharyya coefficient of a query is at most one by $\sqrt{ab}\le\frac{a+b}{2}$ and unit mass,
and it equals one at the dummy. Under world $w$ the mean of $\sqrt{p_{v}/p_{w}}$ at one round is the
coefficient, by the density of the observation law. On a scheduled policy the history law
integrates a product of per-round functions of the observations to the product of their means. The
proof is by induction on the history recursion, through the deterministic query kernel and the
observation kernel. The tournament gives each oriented pair a block of $n$ rounds. The block of a
tested pair asks the pair's best query, which attains the maximum of $H^{2}$ over the finite query
set. The test on a block compares the two likelihoods, with ties broken by a fixed rank of the
worlds, which makes it antisymmetric. Markov's inequality on the square root of the block's
likelihood ratio bounds the error by $\mathrm{BC}^{n}\le e^{-nH^{2}}$, and the combinatorial
tournament law finishes it.

\leanbox{\lean{bc}, \lean{bc\_le\_one}, \lean{bc\_dummy}, \lean{hel}, \lean{bc\_pow\_le},
\lean{lintegral\_sqrt\_ratio}, \lean{lintegral\_prod\_sched}, \lean{block\_error}, \lean{kH},
\lean{qStar}, \lean{tSched}, \lean{beatH}, \lean{beatH\_anti}, \lean{card\_block},
\lean{beatH\_error}, \lean{hellinger\_tournament}, \lean{hellinger\_upper\_star},
\lean{resolved\_le\_scale}, \lean{gChan\_bc}, \lean{hellinger\_tournament\_satisfiable} in
\lean{LeanPaperBHellingerTestV1}}

\paragraph{Shared-query tournament and its limit}
The current upper construction is \lean{share\_tournament} in
\lean{LeanPaperBShareTournamentV1}. Queries are indexed once, with one block of length $n$
per query. Both orientations of a pair choose the same best query (\lean{qS\_symm}),
so the likelihood tests remain antisymmetric even when many pairs reuse observations.
Each individual error bound is still $\mathrm{BC}^n$ (\lean{beatS\_error}); dependence
between tests is harmless for the union bound. The horizon is $|Q|n$, including any dummy
block, hence paid cost is at most that horizon.
The module passed the frozen smoke gate \texttt{smokeSHARP2\_20260925}. The Gaussian private-query
example in \lean{LeanPaperBPrivateQueryV1} is already a concrete instance: its
\lean{private\_query\_lower} forces a cost proportional to the number of private queries,
while \lean{bc\_pq} and \lean{pq\_resolves} instantiate the tournament's information threshold.
It does not prove the old quadratic factor in the number of worlds necessary.

\paragraph{Second-moment obstruction law}
\lean{LeanPaperBChiStarV1} proves the actual history density ratio (\lean{hist\_density}),
the identity $\int p_w^2/p_u=1+\chi^2(p_w,p_u)$ (\lean{m2\_eq\_one\_add\_chi}), and the
history-recursion bound $\mathbb E_w[L_u e^{-tN}]\le1$ (\lean{lintegral\_Z\_exp\_le}).
The unsafe-output events cover the event of insufficient information. Splitting each event
at likelihood ratio $c$ gives $c\delta+e^y/c$ (\lean{one\_alt\_bound}), and optimising $c$
gives \lean{chi\_star\_tail}. The nonnegative-count inequality on its complement yields
\lean{chi\_star\_expect}. Thus the result is about probabilities and expected counts under
the adaptive observation law, not an arithmetic lemma supplied with a transcript bound.
It assumes the stronger finite second moment uniformly over queries; finite KL alone does
not imply that condition. The released proposition is finite-horizon. The stopped KL
proposition elsewhere retains its separately proved scope.

\paragraph{The latent family's rates (Theorem ``Both rates on the latent family'')} The observation
density of a probe is the exact Gaussian mixture over the $8\cdot2^{r}$ latent patterns. Its
characteristic function factors over the latent coordinates, which gives the decision probe's mean
of $\cos Y$ and the tail probe's mean of the parity feature in closed form, with gaps at least
$G_{\rm dec}h$ and $G_{\rm tail}h^{r}$. A nearest-mean test on a bounded feature over the rounds a
schedule gives one query errs with probability at most $2\exp(-(n\,\mathrm{gap}/2)^{2}/(2H))$, by the
concentration bound on the history law. The lower sides are the general law with the per-query
$\chi^{2}$ bounds of the family's channel.

\leanbox{\lean{famLaw}, \lean{charFun\_famLaw}, \lean{head\_factor}, \lean{decMean\_eq},
\lean{tailMean\_eq} in \lean{LeanPaperBFamilyLawV1}; \lean{rawS}, \lean{cntS}, \lean{near\_error},
\lean{dChan}, \lean{safeOut}, \lean{idOut}, \lean{safe\_decoder\_error}, \lean{id\_decoder\_error},
\lean{family\_safe\_minimax}, \lean{family\_id\_minimax} in \lean{LeanPaperBFamilyDecoderV1};
\lean{family\_safe\_stopped\_rate}, \lean{family\_id\_stopped\_rate} in \lean{LeanPaperBStoppedLawV1}}

\paragraph{Tightness on both channels (Theorem ``The general law is tight on both channels'')} For
$t>0$ the pointwise bound
$|p-q|\le\frac t2(\sqrt p-\sqrt q)^{2}+\frac1{2t}(\sqrt p+\sqrt q)^{2}$ integrates to
$\int|p-q|\le tH^{2}+\frac2t$, since $\int(\sqrt p-\sqrt q)^{2}=2H^{2}$ and
$\int(\sqrt p+\sqrt q)^{2}\le4$. The choice $t=4/\int|p-q|$ gives $(\int|p-q|)^{2}\le8H^{2}$. All of this
is non-negative integrals, so no integrability side condition appears. A bounded feature's mean
gap is at most $\int|p-q|$. On the latent family the two probes' gaps then bound the Hellinger
information of every pair below. The decision probe covers pairs that differ in $u$, which every
obstruction of the safe decision contains. The tail probe covers pairs that differ only in $v$. On
the parity channel the separating query, of order $\mathrm{sep}(w,v)=\min(w,v)$ in the paper's
numbering ($\min(w,v)+1$ on the $0$-based indices of the Lean code), has gap at least
$G_{\mathrm{sep}}h^{\mathrm{sep}}\ge G_{m-1}h^{\mathrm{sep}}$, because the gap constant is antitone. So the
Hellinger information dominates the graded information $C_{H}h^{2\,\mathrm{sep}}$, and the proved safe
orders resolve every obstruction at $C_{H}h^{2}$ and $C_{H}h^{2(m-1)}$.

\leanbox{\lean{tvs}, \lean{abs\_sub\_le\_amgm}, \lean{lint\_sqdiff}, \lean{lint\_sqsum},
\lean{tvs\_bound}, \lean{tvs\_sq\_le}, \lean{feature\_gap\_le}, \lean{gap\_sq\_le\_hel},
\lean{fam\_hel\_u}, \lean{fam\_hel\_v}, \lean{fam\_resolves\_safe}, \lean{fam\_resolves\_id} in
\lean{LeanPaperBHellingerFamilyV1}; \lean{Gc\_anti}, \lean{int\_feat\_eq}, \lean{parity\_hel\_sep},
\lean{parity\_kH\_graded}, \lean{resolves\_of\_le} in \lean{LeanPaperBHellingerParityV1}}

\paragraph{The derived safe sets} The family's responses are cube averages. Flipping one
coordinate is an involution of the cube, so a sum whose summand changes sign under some flip is
zero. A coordinate other than $z_{j}$ therefore averages to zero on $\{z_{j}=b\}$, and so does the
tail product when $r\ge2$, since a second tail coordinate is free. The support $K$ has two points,
and the residual of retaining $z_{j}$ at $(t,t,t,1,\dots,1)$ is $(3u-4-c_{j})t+v$, so $v$ cancels in the
oscillation.

\leanbox{\lean{flip}, \lean{sum\_zero\_of\_flip}, \lean{sum\_ind\_coord}, \lean{sum\_ind\_tail},
\lean{resp\_eq}, \lean{F\_Kpt}, \lean{loss\_eq}, \lean{loss\_cases}, \lean{safeOf},
\lean{famSafe\_derived} in \lean{LeanPaperBFamilySafeV1}}

\section{The two layers, the certificate channel and query design}

\paragraph{The decision order} $\mathrm{near}(w,s)$ is $\{w\}\cup\{\lambda:\mathrm{sep}(w,\lambda)>s\}$
and $d(w)$ is the least $s$ at which it has a common safe answer; it exists because every world
has a safe answer and orders are bounded. \emph{Upper side:} at $s=d(w)-1$ there is no common
answer. So every answer safe at $w$ is unsafe at some $\lambda$ with $\mathrm{sep}\ge d(w)$, and
every query sees that $\lambda$ at most $Ch^{2d(w)}$. \emph{Lower side:} a common safe answer $a$
of $\mathrm{near}(w,d(w))$ is unsafe only at worlds of order at most $d(w)$. One query sees each
of them at least $ch^{2\mathrm{sep}}\ge ch^{2d(w)}$, and the uniform allocation loses $1/|Q|$.
Under the ultrametric property, $\mathrm{near}(w,d(w)-1)$ is an obstruction whose pairs all have
order at least $d(w)$.
\leanbox{\lean{near}, \lean{dOrder}, \lean{dOrder\_spec}, \lean{dOrder\_min},
\lean{decision\_upper}, \lean{order\_ge\_obstruction}, \lean{obstruction\_of\_order},
\lean{max\_order\_iff} in \lean{LeanPaperBDecisionOrderV1}; \lean{decision\_lower\_wit} in
\lean{LeanPaperBVisibleV1}}

\paragraph{Blind queries} The earlier form of the lower side, \lean{decision\_lower}, asked for
$ch^{2\mathrm{sq}_q}\le\kappa_q$ at \emph{every} query and pair. On the parity channel a query
that names neither world gives both worlds the same law, so $\kappa_q=0$ there and that
hypothesis fails for $c,h>0$. The endpoint parity results never used it. The repaired form needs
one witness query per pair. A channel whose queries either see a pair, with
$ch^{2\mathrm{ord}}\le\kappa\le Ch^{2\mathrm{ord}}$, or are blind to it, with $\kappa=0$, meets
both hypotheses once every pair is seen: give a blind query an order $N$ above all visible ones.
If the best query of the table is blind, its order is $N$, which bounds the order of the visible
query that exists. On the parity channel query $q\ge1$ sees $(w,w')$ exactly when the masks
differ, which is when $q$ names one of them. Every pair is seen, because at most one of the two
worlds is the unnamed world $m$. The pair's order is $\min(w,w')$ in the paper's numbering.
\leanbox{\lean{decision\_lower\_wit}, \lean{visSq}, \lean{visible\_upper},
\lean{visible\_witness}, \lean{decision\_order\_visible}, \lean{parVis}, \lean{par\_vis\_iff},
\lean{par\_cover}, \lean{parSq}, \lean{par\_sep} in \lean{LeanPaperBVisibleV1}}

\paragraph{The per-query lower bound on the parity channel} Pointwise
$(\sqrt p-\sqrt q)^2\le(p-q)^2/q$, so $2H^2_q\le\chi^2_q$ for every channel and query. At a query
$q$ that sees $(w,w')$ the masks of the two worlds differ at $q$, and the feature $\mathrm{feat}_q$,
bounded by one, has means $c_\pm(q,h,\sigma)$ that differ by at least $G_q(\sigma)h^{q}$
(\lean{gap\_ge}); a mean gap $\Delta$ of a feature bounded by one gives $\Delta^2\le8H^2$. Hence
$\chi^2_q(w,w')\ge\frac14G_q^2h^{2q}\ge\frac14G_{m-1}^2h^{2q}$, the last step because $G_q$ decreases in
$q$. With the upper bound $\chi^2_q\le C_qh^{2q}$ of \lean{parChan\_chi\_le} and $\chi^2_q=0$ at blind
queries, \lean{decision\_order\_visible} applies with order $q$ and $N=m$. For identification the
decision order of world $w$ (paper numbering) is $\min(w,m-1)$: below it the world $m$, or the
world $m-1$ when $w=m$, is still inseparable from $w$. An instance ($m=3$, $\sigma=1$, $h=\frac12$)
shows the hypotheses are met.
\leanbox{\lean{sqrt\_sub\_sq\_le}, \lean{hel\_le\_chi}, \lean{par\_gap\_visible}, \lean{par\_chi\_visible},
\lean{cW\_pos}, \lean{par\_chi\_visible\_unif}, \lean{parKap\_upper}, \lean{parKap\_blind}, \lean{parKap\_lower},
\lean{parity\_world\_law}, \lean{parity\_id\_order}, \lean{parity\_id\_world\_law}, \lean{parity\_world\_instance} in
\lean{LeanPaperBParityWorldV1}}

\paragraph{The parity cancellation} Expand $(\sum_\ell z_\ell)^k$ as a sum over maps
$f:[k]\to[j]$ of $\prod_iz_{f(i)}$. Summing over the cube factorises by coordinates:
$\sum_z\prod_\ell z_\ell^{1+|f^{-1}(\ell)|}=\prod_\ell\bigl(1+(-1)^{1+|f^{-1}(\ell)|}\bigr)$. A
factor vanishes when some fibre is empty. For $k<j$ every $f$ misses a coordinate. For $k=j$ only
the surjections survive, each fibre has one element, each contributes $2^j$, and there are $j!$
of them.
\leanbox{\lean{sg}, \lean{pow\_sum\_expand}, \lean{prod\_fiber}, \lean{term\_eq},
\lean{moment\_expand}, \lean{factor\_zero}, \lean{parity\_moment\_lt}, \lean{card\_surj},
\lean{parity\_moment\_eq} in \lean{LeanPaperBVisibleV1}}

\paragraph{The constants, numerically} \NUM{} In the expansion of the query-$j$ density in $h$,
the $k$-th term averages $(\sum_\ell Z_\ell)^k\varphi_\sigma^{(k)}(y\mp\prod_\ell Z_\ell)$. Writing
$\varphi_\sigma^{(k)}(y\mp s)$ for $s=\pm1$ as the average of its two values plus $s$ times half
their difference, only the second part depends on the mask, and the cancellation leaves
$p_+-p_-=h^ja_j+O(h^{j+1})$ with
$a_j(y)=(-1)^jj^{-j}(\varphi_\sigma^{(j)}(y-1)-\varphi_\sigma^{(j)}(y+1))$. This step and its
remainder are not formalized. With $p_0=\frac12(\varphi_\sigma(\cdot-1)+\varphi_\sigma(\cdot+1))$ and
$c_j=\frac12\int a_j^2/p_0$, the prediction is $\mathrm{KL}\sim c_jh^{2j}$ in both directions and
$H^2\sim c_jh^{2j}/4$. Quadrature (\texttt{paperB\_parity\_constants.json}, $4/4$ checks, among
them the cube identity above by exact enumeration) gives:
\begin{center}
\footnotesize
\begin{tabular}{@{}ccc@{}}
\toprule
$j$ & $c_j$, $\sigma=1$ & $c_j$, $\sigma=\frac12$\\
\midrule
$1$ & $1.468$ & $7.827$\\
$2$ & $0.1862$ & $2.455$\\
$3$ & $1.172\times10^{-2}$ & $0.6558$\\
$4$ & $4.794\times10^{-4}$ & $0.1252$\\
$5$ & $1.472\times10^{-5}$ & $1.752\times10^{-2}$\\
\bottomrule
\end{tabular}
\end{center}
At $h=0.025$ every ratio $\mathrm{KL}/(c_jh^{2j})$ and $4H^2/(c_jh^{2j})$ lies within $0.6\%$ of
$1$, and the distance shrinks with $h$. Fed to the critical game with the visibility table, these
give the identification constant at world $w$. For $w\le m-1$ the critical alternatives are the
worlds $\lambda>w$, which only query $w$ sees at order $w$, so $G^*=c_w$. For $w=m$ the only
critical alternative is $m-1$, so $G^*=c_{m-1}$. This instance of Theorem ``The limit of the
rescaled game'' is derived here and not formalized.

\paragraph{The critical limit} For an allocation $\alpha$ the inner value is a finite minimum, and
the game value is its supremum over the simplex, bounded by a single alternative's row.
\begin{itemize}
\item \emph{Upper:} eventually $|f_h-\hat c|<\varepsilon/4$ on the critical alternatives, and
$\sum_q\alpha_q=1$ makes the bound uniform in $\alpha$.
\item \emph{Lower:} take a critical allocation within $\varepsilon/4$ of $G^*$ and mix in the
uniform one with weight $\beta=\min\{\frac12,\varepsilon/(4(G^*+1))\}$. The critical rows move by
at most $\beta G^*$, and every fast row eventually exceeds $G^*$ through its diverging query,
which now has weight at least $\beta/|Q|$.
\end{itemize}
The expansion lemma splits $h^{2\mathrm{sq}}=h^{2s}h^{2(\mathrm{sq}-s)}$ on critical rows and
$h^{2s}=h^{2\mathrm{sq}}h^{2(s-\mathrm{sq})}$ on fast rows.
\leanbox{\lean{inner}, \lean{gval}, \lean{mixU}, \lean{critical\_limit}, \lean{pow\_tendsto\_zero},
\lean{inv\_pow\_tendsto\_atTop}, \lean{expansion\_hyps}, \lean{decision\_limit},
\lean{inner\_smul}, \lean{gval\_smul}, \lean{private\_game}, \lean{shared\_game},
\lean{same\_order\_ratio} in \lean{LeanPaperBCriticalGameV1}}

\paragraph{The certificate channel} The observation's region is a clamped floor. Inside a region
the density is the same restricted Gaussian in every world, so the per-query KL is a finite sum
over the $k$ symbols. For $\delta\le1/k$ that sum is at most $\log(1/((k-1)\delta))/(k-1)$.
\leanbox{\lean{reg}, \lean{massR}, \lean{klq\_cert\_one}, \lean{kl\_sum\_eq}, \lean{klq\_cert\_le},
\lean{certSafe\_unsafe}, \lean{cert\_error}, \lean{cert\_cost}, \lean{cert\_tight} in
\lean{LeanPaperBCertificateChannelV1}}

\paragraph{Query design} The hitting form is a rewriting of $\mathrm{sep}_M$ as a minimum over
the menu. In the embedding, an obstruction contains two distinct worlds, since each world's
only safe answer is itself. If they belong to different pairs every query separates them at
order $0$. Otherwise they are twins, and only a query of their set separates them at order $1$.
\leanbox{\lean{sepM}, \lean{DecLE}, \lean{IdLE}, \lean{dec\_hitting}, \lean{id\_implies\_dec},
\lean{dec\_mono}, \lean{gap\_example} in \lean{LeanPaperBDecisionOrderV1}; \lean{hsSafe},
\lean{hsSq}, \lean{twins}, \lean{twins\_unsafe}, \lean{twins\_pair}, \lean{hitting\_embeds},
\lean{order\_zero\_fails}, \lean{order\_two}, \lean{id\_iff\_hitting} in
\lean{LeanPaperBQueryDesignV1}}

\section*{Build record}
The corpus is verified by \texttt{verify\_all.sh}, which freezes its inputs, builds every listed
module and every local import from an empty cache in import order, counts \texttt{sorry} and
\texttt{native\_decide} outside comments, and checks every audited axiom closure against
\{\texttt{propext}, \texttt{Classical.choice}, \texttt{Quot.sound}\}. Every module cited in the
paper and in this supplement, $52$ in all, passed a gate whose frozen copy is byte-identical, by
hash, to the current source.
\begin{itemize}
\item The last full gate, \texttt{full7\_20260925} ($219$ modules, $1704$ theorems), covers $40$ of
them, including \lean{LeanUnifiedObstructionOrderV1}, \lean{LeanUnifiedQueryOrderV1},
\lean{LeanPaperBHellingerTestV1} and \lean{LeanPaperBStoppedLawV1}.
\item \texttt{smokeR2\_20260925} ($96$ modules, $771$ theorems): \lean{LeanPaperBFamilySafeV1},
\lean{LeanPaperBHellingerFamilyV1}, \lean{LeanPaperBHellingerParityV1}.
\item \texttt{smokeUP1\_20260925} ($79$ modules, $762$ theorems):
\lean{LeanPaperBCertificateChannelV1}, \lean{LeanPaperBChiStarV1}, \lean{LeanPaperBPrivateQueryV1},
\lean{LeanPaperBShareTournamentV1}.
\item \texttt{smokeUP2\_20260925} ($107$ modules, $934$ theorems): \lean{LeanPaperBCriticalGameV1},
\lean{LeanPaperBQueryDesignV1}.
\item \texttt{smokeUP4\_20260925} ($139$ modules, $1199$ theorems): \lean{LeanPaperBDecisionOrderV1},
\lean{LeanPaperBVisibleV1}.
\item \texttt{smokeB1\_20260926} ($44$ modules, $353$ theorems): \lean{LeanPaperBParityWorldV1}.
\end{itemize}
Each smoke gate freezes its roots with their full import closure and applies the same
acceptance conditions: exit status $0$, no \texttt{sorry}, no \texttt{native\_decide}, standard
axioms, no drift.

\begin{thebibliography}{00}

\bibitem{duc2015} A. Duc, S. Faust, and F.-X. Standaert, ``Making masking security
proofs concrete,'' IACR ePrint 2015/119, 2015.

\bibitem{poli2023} M. Poli et al., ``Ideal abstractions for decision-focused learning,''
arXiv:2303.17062, 2023.


\bibitem{javdani2014} S. Javdani, Y. Chen, A. Karbasi, A. Krause, D. Bagnell, and
S. Srinivasa, ``Near optimal Bayesian active learning for decision making,'' in
\emph{Proc. 17th Int. Conf. Artificial Intelligence and Statistics (AISTATS)}, 2014,
pp.~430--438.

\bibitem{grosso2014} V. Grosso, F.-X. Standaert, and E. Prouff, ``Low entropy masking
schemes, revisited,'' in \emph{Smart Card Research and Advanced Applications (CARDIS)},
2013, pp.~33--43.

\bibitem{rioul2021} O. Rioul, W. Cheng, and S. Guilley, ``Cumulant expansion of mutual
information for quantifying leakage of a protected secret,'' in \emph{Proc. IEEE Int. Symp.
Information Theory (ISIT)}, 2021, pp.~2596--2601.

\bibitem{degenne2019} R. Degenne and W. M. Koolen, ``Pure exploration with multiple correct
answers,'' in \emph{Advances in Neural Information Processing Systems 32}, 2019,
pp.~14591--14600.

\bibitem{wangyou2026} Q. Wang and W. You, ``Ranking-and-selection with multiple correct answers and non-answerable estimates,'' arXiv:2606.21889, 2026.
%% previous (wrong) entry kept for the record: Z. Wang and K. You, ``Information games with multiple correct answers
%% and piecewise-convex alternative sets,'' arXiv:2606.21889, 2026.

\bibitem{odonnell2014} R. O'Donnell, \emph{Analysis of Boolean Functions}. Cambridge,
U.K.: Cambridge Univ. Press, 2014.

\bibitem{tsybakov2009} A. B. Tsybakov, \emph{Introduction to Nonparametric Estimation}.
New York, NY: Springer, 2009.

\bibitem{mathlib2020} The mathlib Community, ``The Lean mathematical library,'' in
\emph{Proc. 9th ACM SIGPLAN Int. Conf. Certified Programs and Proofs (CPP)}, 2020,
pp.~367--381.

\bibitem{demoura2021} L. de Moura and S. Ullrich, ``The Lean 4 theorem prover and
programming language,'' in \emph{Proc. 28th Int. Conf. Automated Deduction (CADE)}, 2021,
pp.~625--635.

\end{thebibliography}

\end{document}
